\section{Dispersion for abstract sources}\label{sec:pi-psdt}

This section proves a local implication for an abstract source.
It uses Linnik's dispersion method \cite{Linnik63}. The second-dispersion
setup has close precedents in Maynard \cite{MaynardII} and Lichtman
\cite[\S5]{Lichtman23}; the proof below tracks its own coefficient
masks, restoration terms, and moving residue.
The local hypotheses are (H1)--(H5). The source-wide estimate (H6)
and the box-family coefficient bound (H7) are used only in the
global summation in Section~\ref{sec:pi-global}.
Each application verifies the required conditions for its own measures.
The geometric and source partition results appear below.
Theorem 2.8 is stated in
Section 2, (2.38)--(2.43), and derived from the classical inputs in
\S\ref{p1:sec:theorem-2-8-proof}. We keep its full moving-frequency bracket.
The local estimate has three terms.
The restoration error $E_{\rm GR}$ is bounded apart from the spectral term.

\subsection{Notation and source hypotheses}

\emph{Local notation.} Put $L=\log X$; \S\ref{sec:pi-global} writes $\log N$ explicitly. $T$ is the lower endpoint of the selected Type-II cell, as in Lemma~\ref{lem:aa37-typeII-lengths} for the prime source. Here $\vartheta_{\rm sp}=7/32$. The zero-mode functional is $\mathcal Z(V)$ of (3.31); $Z$ is the scale of $z$. $B$ is a Heath--Brown box. $D$ and $C$ are the scales of $b$ and $c$ in (3.62), as in Theorem 2.8. $B_{z,y}$ is the count (3.70). $S$ is the scale of $\gamma_2$. $S_B^{\rm PSDT}$ denotes the signed box sum in (3.12). Type I-1 and I-2 denote boxes with one and two large variables (\S\ref{sec:pi-nonII}). Since $\log X=\log N+O(\log\log N)$, fixed powers of $\log X$ and $\log N$ are comparable.
$I_T$ is the selected integer interval of a Type-II cell; $I_N$ is reserved for the source interval of Theorem~\ref{thm:aa37-interior}.
Some letters have local meanings, fixed at each use.
The letter $C$ is a fixed logarithmic exponent in (3.13), (3.92) and similar bounds, and the scale of $c$ from (3.62) on.
Similarly, $\varepsilon$ is the source cut of Theorem~\ref{thm:aa37-interior}, $\epsilon$ is a small fixed exponent, and $\varepsilon_y\in\{\pm1\}$ is a sign.
The letter $Y$ is a summation length in (3.16)--(3.19) and (3.35).
It is the auxiliary scale $H_*SF$ in the spectral step near (3.70).
In Theorem~\ref{thm:aa37-interior} it is the smooth-prime cutoff $D^{\sigma^2}$.
$\mathcal T$ is the tuple set of \S\ref{sec:pi-psdt}; it is unrelated to $T$.

Let $N$ be a sufficiently large even integer and let $a$ be a positive
integer. Fix $0<K_{\rm scale}<\infty$ and

\[
0<\sigma\le10^{-10},\quad \tau=\sigma^9,\quad\delta=30\sigma,
\quad\vartheta_{\rm sp}=7/32,\quad\eta=\sigma/1000,
\]
\[
\eta_L=10^{-4},\qquad
\ell_0=\frac{53}{412}-\frac{\eta_L}{2}
       =\frac{264897}{2060000},\qquad
c_\sigma=\ell_0\sigma^2/4.
\tag{3.1}
\]

The local product scale is
$N/(\log N)^{K_{\rm scale}}\le X\le N$.
All implied constants in this section may depend on the fixed
$K_{\rm scale}$ and the fixed source bounds. They may also depend on
the parameters displayed at the point of use. Put
\[
\begin{gathered}
L=\log X,\quad z_0=X^{1/(\log\log X)^3},\quad
u_0=N^{\ell_0\sigma^2},\\
p_0=\min\{p\text{ prime}:p\nmid 2aN\}.
\end{gathered}
\tag{3.2}
\]
In particular, $N/u_0=N^{1-\ell_0\sigma^2}$.

Thus $p_0\ge3$ and $p_0\nmid aN$.
Every prime below $p_0$ divides $2aN$.
Chebyshev's bound $\sum_{p<p_0}\log p\gg p_0$ gives
$p_0\ll\log(2aN)$. By C4 below, $a\le X(\log N)^{K_{\rm scale}}$,
so $p_0\ll\log N<z_0\le u_0$ eventually. Write

\[
\mathcal S_t(q;a)=1_{t\equiv a\ (q)}-
                    \frac{1_{(t,q)=1}}{\varphi(q)}.
\tag{3.3}
\]

A \textbf{Class-II box} consists of the following data.
These are source hypotheses.

\textbf{C1. Exact unrestricted coefficient.} Fix a source term, a prefix cell,
a sign of the well-factorable remainder in Iwaniec's theorem \cite[Theorem~4]{AA37},
and a tail-label vector $b$. Its coefficient is exactly

\[
C_b(q)=\sum_{d,(p_j)}\lambda_d
                  1_{q=d\prod_{j\in J}p_j},\qquad |\lambda_d|\le1,
\tag{3.4}
\]

where each $p_j$ ranges independently over its fixed prime set and
$p_j\ge u_L$.
Here $u_L\ge u_0$ is the source's lower cutoff for labelled tail primes,
and $P(u_L;N)=\prod_{p<u_L,\,p\nmid2N}p$.
The variable $d\mid P(u_L;N)$ is squarefree with the
cap of that theorem. The full modulus is at most $X^{0.6}$.
Every supported full modulus is prime to $2N$.
Selectors are constant on the block.
Exact identities remove the equal-tail, prefix-diagonal, natural-comparator,
cut, and selector-transfer corrections.
Each is a separate contribution outside (3.4).

\textbf{C2. Exact factors.} There is a labelled partition
$J=J_1\sqcup J_2\sqcup J_3$, with $d$ attached to $J_1$, such that

\[
C_b=C_{\rm slot}(\gamma_1*\gamma_2*\gamma_3),\qquad
C_{\rm slot}=\prod_{i=1}^3|J_i|!=O_\sigma(1),\qquad |\gamma_i|\le1.
\tag{3.5}
\]

The last two factors are $u_0$-rough. If $J_s=\varnothing$ for $s=2,3$, then $\gamma_s(m)=\mathbf1_{m=1}$. The factor $\gamma_1$ always contains the smooth-divisor sum of Iwaniec's theorem. If $J_1=\varnothing$, then $\gamma_1(m)=\sum_d\lambda_d\mathbf1_{m=d}$.
Equality (3.5) holds on all supported moduli,
with repeated-prime multiplicities.
The source's labelled-tuple map proves it.
It covers moduli with repeated primes as well as squarefree ones.

\textbf{C3. Packing.} Localize dyadically the first factor $\gamma_1$
at $\asymp Q$, $\gamma_3$ at $\asymp R$, and $\gamma_2$
at $\asymp S$. Here $T$ is the lower endpoint of the selected Type-II
cell; $N_I$ is its integer length.
Enlarge the sufficiently-large threshold for $N$ to absorb fixed dyadic constants. Then

\[
Q<TX^{-29\sigma},\qquad QRS<M_I X^{-89\sigma},
\qquad N_I\ge QX^{\sigma/3},\quad M_I\ge X^{1/2}.
\tag{3.6}
\]

Here $Q$ is the scale of the first factor $\gamma_1$ (the outer
slot scale, written $Q_{\rm slot}$ in Appendix~\ref{app:proof-aa37}).
It is not the scale of the full modulus $q$ of (3.4), which is $QRS$.
Thus (3.6) gives the first inequality in (3.7).
We list the four spectral packing inequalities together.
\[
Q<TX^{-\sigma},\quad T^2R^2S<X^{1-8\sigma},
\quad A_0^{\vartheta_{\rm sp}} TQR^2S^5<X^{2-15\sigma},
\]
\[
(TR/S)^{2\vartheta_{\rm sp}}TQR^2S^5<X^{2-15\sigma},
\qquad A_0=X^{1+\sigma/20}.
\tag{3.7}
\]
Each application must verify the stated packing and interval
conditions for its own source. Appendix~\ref{app:proof-aa37} does so
for the prime source in Theorem~\ref{thm:aa37-interior}.

\textbf{Removed squareful and other restricted terms.}
A source may remove forbidden tails, forbidden prefixes or
retained allowed repeats.
Their contributions are bounded separately.
One admissible bound is $O_\sigma(N^{1-\ell_0\sigma^2}L^4)$,
after all source-label and representation multiplicities.
Theorem~\ref{thm:main} concerns the unrestricted coefficient (3.4).
The removed terms need their own estimates.

The prime source in Theorem~\ref{thm:aa37-interior} is
$\mathbf1_{\mathbb P}(n)\mathbf1_{N/(\log N)^B<n<(1-\varepsilon)N}$.
Appendix~\ref{app:proof-aa37} verifies its coefficient, interval,
cutoff, support, and boundary conditions. The hypotheses above also
state the local estimate for any other source that satisfies them.

Apply the finite HB expansion below.
Then split into disjoint short product boxes of
relative width $L^{-C_{\rm box}}$. Boundary boxes contribute separately.
On each interior Class-II box select an HB variable in the window.
Maximize its lower-endpoint exponent $e_i=\log_N V_i$.
Break ties by the fixed variable index, as in the classification in
Appendix~\ref{app:proof-aa37}.
Call it $n$, and the product of the other variables $m$.

In (3.8), allow a roughness cutoff $\zeta\ge z_0$,
where $z_0=z_0(X)$ is the baseline in (3.2).
All estimates below are uniform in $\zeta$.
Upper bounds use only $P^-(n)\ge z_0$.
The SW proof (3.25) for the selected factor uses the cutoff $\zeta$.
It remains uniform under stronger roughness.

\[
\alpha_n\in\left\{\mu(n),1,\frac{\log n}{L}\right\}
1_{P^-(n)\ge \zeta}1_{I_T}(n),\quad |\alpha_n|\le1,
\qquad |\beta_m|\ll\tau_{23}(m).
\tag{3.8}
\]

The fixed HB binomial coefficients and tags may be kept as exterior
scalars. $I_T\subset(T,2T]$ has length $N_I\asymp TL^{-C_{\rm box}}$.
The $m$-interval is majorized by

\[
\psi_I(m)=\phi_0((m-m_0)/M_I),\qquad
M_I\asymp (X/T)L^{-C_{\rm box}},
\tag{3.9}
\]

Fix $\phi_0\in C_c^\infty([-1,1])$, with $0\le\phi_0\le1$ and $\phi_0=1$ on $[-1/2,1/2]$.
It is independent of $N$, $X$, the source and the box.
It is also independent of all shell and tuple variables.
For every $j\ge0$, the seminorms $\|\phi_0^{(j)}\|_1$ and $\|\phi_0^{(j)}\|_\infty$ are fixed constants.
The support of $\psi_I(m)=\phi_0((m-m_0)/M_I)$ lies in the doubled interval.
The majorant depends on the box only through $m_0$ and $M_I$.
On that entire support the cut geometry gives

\[
1\le |mn-a|\le X^3.
\tag{3.10}
\]

\textbf{C4. Moving frequency.} All original moduli are prime to $a$,
and $0<a\le X(\log N)^{K_{\rm scale}}\le A_0$. This is the moving
frequency in the spectral theorem. It may vary with $N$.

\textbf{C5. Window.} The source selection gives

\[
N^{60\sigma}\le T\le N^{1/3+14\sigma},\qquad
X^{60\sigma}\le T\le X^{1/3+15\sigma}.
\tag{3.11}
\]

The second assertion follows from the first and
$\log(N/X)/\log N=o(1)$, with $N_0$ chosen after fixed
$\sigma$. \textbf{Claim P (nonresonance):} the doubled support satisfies (3.10).
The source assumptions for the non-II
extension are in \S\ref{sec:pi-nonII}.

\subsection{Statement and theorem inputs}

\textbf{Length hypothesis for the selected integer interval (LEN).}
Uniformly over the selected boxes, the interval has length
$N_I\asymp T L^{-C_{\rm box}}$ and satisfies
$N_I\ge(2T)^{1-2^{-126}}$.
The integer interval used in (3.17) has this same length, with the
endpoint convention specified by the source.
Lemma~\ref{lem:aa37-typeII-lengths}, (B.22)--(B.27), proves
this for the prime source.
Each application proves it for its own source.

For the coefficient of one Heath--Brown box define

\[
\Delta_B(q;a)=\sum_{\mathbf v\in B}w(\mathbf v)
                 \mathcal S_{\prod v_i}(q;a),\qquad
S_B^{\rm PSDT}=\sum_q C_B(q)\Delta_B(q;a).
\tag{3.12}
\]

\subsection{An abstract dispersion estimate}

The local argument is source-independent.
We state it before the verification for specific sources.
Let $\mathcal B$ be a finite
family of dyadic boxes at scale $X$, and let $a$ be a moving frequency.
For $B\in\mathcal B$ write $C_B(q)$ for its signed modulus coefficient.
Define $\Delta_B(q;a)$ by (3.12).

The following hypotheses define the class used in this paper.

\begin{description}[leftmargin=2.4em,style=nextline]
\item[(H1) Exact coefficient] On every box, $C_B(q)$ has the exact
form (3.4), with its fixed labelled prime sets, squarefree
smooth-divisor variable $d$, and repeated-prime multiplicities. Every
supported modulus is coprime to $2a$.

\item[(H2) Three-factor decomposition] The labelled variables split
into three groups satisfying (3.5). The last two factors are
$u_0$-rough, as in C2. The identity holds on all
supported moduli, including repeated-prime multiplicities.

\item[(H3) Geometric packing and box data]  After dyadic localization,
the scales $Q,R,S$ and the test scale $T$ satisfy (3.6)--(3.7).
On each box, the Heath--Brown coefficients factor as $\alpha_n\beta_m$
with (3.8).  The smooth majorant has the form (3.9).
Its center satisfies $|m_0|\le C_{\rm ctr}X/T$, where $C_{\rm ctr}$
is a fixed source-class constant, uniform over all boxes, shells,
tuple variables, and permitted $N,X,a$.
Its full doubled support satisfies (3.10).
The selected integer interval satisfies LEN, including its strong
length bound and integer endpoint convention.
For this abstract dispersion estimate, assume modulus support in
$q\le X^{\vartheta_0}$, with $\vartheta_0$ as in (3.12a).

\item[(H4) Moving frequency]  The frequency is the positive integer $a$,
with $N$ even, $0<a\le X(\log N)^{K_{\rm scale}}\le A_0$, and all original moduli are coprime
to $a$.

\item[(H5) Selected window]  The tested scale obeys (3.11), with all
parameters fixed after $\sigma>0$ is fixed.

\item[(H6) Source-wide remainders and boundary] Let a source supply
a logarithm-freezing estimate and boundary estimate for the full box family.
After summing over all source labels, shells, boxes, and representation
multiplicities, these estimates must give
\[
E_{\rm log}\ll N(\log N)^{c_{\log}-C_{\rm box}},\qquad
E_{\rm bdry}\ll N(\log N)^{c_R-C_{\rm box}/2},
\]
with fixed $c_{\log}$ and $c_R$. This condition is used in the
global summation, not in the local theorems below.

\item[(H7) Box-family coefficient bound] The finite box family
$\mathcal B=\{B_j\}$ carries outer coefficients $c_j$, coming from
the source labels and the Heath--Brown identity. At fixed $\sigma$
they satisfy (3.91):
$\sum_j|c_j|\le C_{\sigma,C_{\rm box},\mathcal C_{\rm src}}(\log N)^{K_0}$ with a fixed
$K_0\le24C_{\rm box}+26$. This condition is used only for the
aggregate bound (3.92), not for the per-box estimates.
\end{description}

Write $\mathcal C_{\rm src}$ for the fixed source constants in
(H1)--(H5), such as $C_{\rm slot}$, $C_{\rm ctr}$ and the implied
constants in (3.8)--(3.11).

The Siegel--Walfisz bound (3.25) for the selected factor follows in
\S\ref{sec:pi-sw} from the classical theorem \cite[Corollary~5.29]{IK04}. BFI \cite{BFI86} and
Theorem 2.8, proved from classical inputs in
\S\ref{p1:sec:theorem-2-8-proof}, are theorem inputs, not source hypotheses.
The source-wide boundary bound (3.24) is hypothesis (H6). It is used
in the source estimates, not in the per-box theorem (3.13).

Put
\[
 \Theta_0=\frac{19101}{32000},\qquad
 \vartheta_0=\Theta_0-\eta_L>\frac12.
\tag{3.12a}
\]
The number $\Theta_0$ is the small-atom empty-prefix packing level
of Proposition~\ref{prop:pack-ceiling} in \S\ref{sec:pack}.
The exponent $\vartheta_0$ is a support convention in the abstract
dispersion estimate. The actual support level comes from the packing inequalities
(3.6)--(3.7) after all Iwaniec factors are restored.
In Theorem~\ref{thm:aa37-interior}, (B.2)--(B.4) give
$q<N^{e_Q}$ with $e_Q<\Theta(u)\le\Theta(0)=153/256<0.6$; hence
$q\le X^{0.6}$ on every source shell for sufficiently large $N$.
The fixed exponent $\vartheta_0$ belongs to the abstract support
hypothesis of the introductory source-box estimate. The variable ceiling
$\Theta(u)=(5-7/32-3u)/8$ in Theorem~\ref{thm:aa37-interior}
controls the Iwaniec coefficient family as $u$ varies.
At $u=1/500$, $\Theta(u)=19101/32000=\Theta_0$,
whereas $\vartheta_0=\Theta_0-10^{-4}$.
The introductory estimate applies to its stated abstract support class;
the Iwaniec result at this endpoint uses the $X^{0.6}$ theorem below.
Theorem~\ref{thm:standalone} is the abstract dispersion estimate under
the support and packing hypotheses in (H3), including the stronger
support $q\le X^{\vartheta_0}$.
Theorem~\ref{thm:main} is the dispersion engine used for the actual
source support. Theorem~\ref{thm:aa37-interior} in the application
uses that engine together with the source estimates of
Appendix~\ref{app:proof-aa37}. At $u=1/500$ it gives the limiting
ceiling $\Theta(1/500)$ for this specified signed mean value.

\begin{theorem}[Dispersion estimate under the (H3) packing hypotheses]\label{thm:standalone}
\textup{(a)} Assume (H1)--(H5). Then, for every fixed $A'>0$ and
every box $B$,
\[
 \left|\sum_{q\le X^{\vartheta_0}}C_B(q)\Delta_B(q;a)\right|
 \ll_{\sigma,K_{\rm scale},C_{\rm box},A',\mathcal C_{\rm src}}
 XL^{-A'}+X^{1-c_\sigma}L^C,
\tag{3.12b}
\]
where $c_\sigma=\ell_0\sigma^2/4>0$ and
$C=C(\sigma,C_{\rm box},A',K_{\rm scale},\mathcal C_{\rm src})$ is the exponent of (3.13).
The implied constant and the threshold for $X$ are uniform in the box,
the even integer $N$, and all moving frequencies $a$ permitted by
(H1)--(H5). They may
depend on the fixed source bounds and on $\sigma$. They are
ineffective, since the proof uses the Siegel--Walfisz theorem in (3.25).
Under (H3) every supported modulus satisfies $q\le X^{\vartheta_0}$, so
the truncation in (3.12b) is vacuous and its left side is $|S_B^{\rm PSDT}|$.

\textup{(b)} Assume in addition (H7). Fix $A>0$ and take
$A'=A+24C_{\rm box}+31$ in \textup{(a)}. Then the weighted sum
$E_{\rm II}=\sum_j|c_j||S_{B_j}^{\rm PSDT}|$ over the box family
satisfies (3.92), that is,
$E_{\rm II}\ll_{\sigma,A,C_{\rm box},K_{\rm scale},\mathcal C_{\rm src}}
N(\log N)^{-A}$.
This bound is uniform in the scales $X$ of the boxes. For
$N/(\log N)^{K_{\rm scale}}\le X\le N$ one has
$\log X\ge\log N-K_{\rm scale}\log\log N\ge\tfrac12\log N$ for large $N$,
so $X(\log X)^{-A'}\le2^{A'}N(\log N)^{-A'}$ and
$X^{1-c_\sigma}(\log X)^C\le N^{1-c_\sigma}(\log N)^C$.
Part \textup{(b)} holds for any fixed $C_{\rm box}>0$. The condition
$C_{\rm box}\ge\lceil2A+32\rceil$ of Theorem~\ref{thm:aa37-interior}
serves only the source-wide errors of (H6).

Under the support hypothesis in (H3), the theorem gives arbitrary
logarithmic saving for these signed sums through
$q\le X^{\vartheta_0}=X^{0.59680625}$, where $\vartheta_0>1/2$.
\end{theorem}

\begin{proof}
The support bound assumed in (H3) makes the sum in (3.12b)
the full sum $S_B^{\rm PSDT}$ of (3.12).
Hypotheses (H1)--(H5) imply those of Theorem~\ref{thm:main}.
Hence Theorem~\ref{thm:main} applies to $B$ and gives (3.13), which is (3.12b).
This proves \textup{(a)}. Part \textup{(b)} follows from \textup{(a)}
and (H7) by the summation (3.92) in \S\ref{sec:pi-global}.
\end{proof}

\begin{theorem}[Dispersion theorem]\label{thm:main}
Assume (H1)--(H5), where in (H3) the support bound $q\le X^{\vartheta_0}$
may be relaxed to $q\le X^{0.6}$.
Thus each box satisfies (3.8)--(3.10) and LEN.
Then, for each fixed $A'>0$,

\[
|S_B^{\rm PSDT}|\ll_{\sigma,K_{\rm scale},C_{\rm box},A',\mathcal C_{\rm src}}
 XL^{-A'}+X^{1-c_\sigma}L^C.
\tag{3.13}
\]

Here $C=C(\sigma,C_{\rm box},A',K_{\rm scale},\mathcal C_{\rm src})$.
It multiplies only power-saving terms.
It enters only the final choice of $N_0$.
The exponent $C$ is fixed after the proof parameters.
It is uniform in the box, each class in the random partition, the residue,
and the moving shift.
Constants and $N_0$ may depend on the fixed $\sigma$.
All bounds are for fixed $\sigma>0$ as $N\to\infty$.
The implied constants and the threshold $N_0$ are ineffective, since
the proof uses the Siegel--Walfisz theorem in (3.25).
The estimate (3.13) is per box and does not use (H6) or (H7).
If (H7) also holds and $A'=A+24C_{\rm box}+31$, then the weighted sum
over the box family satisfies (3.92).
\end{theorem}

Hypotheses (H1)--(H5) give the coefficient, factor identity, packing, box,
frequency and window conditions, including LEN and the integer endpoint convention.
BFI \cite[Theorem~0(a)]{BFI86} supplies (3.48), Theorem 2.8 of
\S\ref{p1:sec:theorem-2-8-proof} supplies (3.14)--(3.15), and
(3.25) supplies SW.
The source-wide condition (H6) enters the source error terms, including
(B.29) in the proof of Theorem~\ref{thm:aa37-interior} and the auxiliary
bound (3.95). It is not used in the per-box Theorem~\ref{thm:main}.
The box-family bound (H7) enters only (3.92).
The spectral part gives a term $X^{1-1.725\sigma}L^C$ (\S\ref{sec:pi-spectral}).
It is absorbed in $X^{1-c_\sigma}L^C$, since $c_\sigma<\sigma/2<1.725\sigma$.

The dependencies have the following order:
\[
\begin{gathered}
\text{source identities, box geometry and source estimates}\ \Longrightarrow\ (\mathrm{H1})\text{--}(\mathrm{H7}),\\
\text{Section 2 and (3.25)}\ \Longrightarrow\ \text{BFI, Theorem 2.8 and SW},\\
\text{(H1)--(H5)}+\text{theorem inputs}\ \Longrightarrow\ (3.13),\\
(3.13)+(\mathrm{H7})\ \Longrightarrow\ (3.92).
\end{gathered}
\]
The hypotheses are stated independently of any particular source.
Appendix~\ref{app:proof-aa37} verifies them for
Theorem~\ref{thm:aa37-interior}.
Source comparison and the estimates for removed squareful terms are separate steps.
Aggregate logarithm-freezing estimates enter in \S\ref{sec:pi-global}.
Theorem~\ref{thm:main} uses (H1)--(H5) and the stated theorem inputs.

The inputs are the following theorems, printed in Section 2.
The Siegel--Walfisz theorem (SW) \cite[Corollary~5.29]{IK04} gives prime distribution in reduced progressions
to moduli at most a fixed power of the logarithm.
\cite[Theorem~0(a), (2.2)]{BFI86}, with \cite[(A2)]{BFI86},
gives the BDH quadratic estimate.
It applies to sequences satisfying its norm-normalized SW condition.
On the domain $(rd,sc)=1$, Theorem 2.8 gives

\[
\begin{aligned}
&\left|\sum_{(r,s)=1}\sum_{0<n\le N_T}b_{n,r,s}
 \sum_{c,d:(rd,sc)=1}g(c,d,n,r,s)
 e\left(\frac{a_Tn\overline{rd}}{sc}\right)\right|\\
&\qquad\ll (2+A_T C D N_T R_T S_T)^{\epsilon+18\epsilon_0}J_*(a_T).
\end{aligned}
\tag{3.14}
\]

Here $0<|a_T|\le A_T$, where $A_T$ is the frequency bound $A$ of Theorem 2.8, and $b=0$ unless $(a_T,rs)=1$. The positive
scales meet the endpoint conditions of Section 2. The bracket is

\[
\begin{aligned}
J_*^2={}&\sum_{n''\mid a_T^\infty}A_T^{\vartheta_{\rm sp}}(n'')^{\vartheta_{\rm sp}}
 \bigl[C S_T(N_T/n''+R_TS_T)(C+DR_T)\\
&\qquad+|a_T|N_TR_TS_T\bigr]\|b^\flat(n'')\|_2^2\\
&+\bigl[C S_T(CD\sqrt{R_T})^{2\vartheta_{\rm sp}}
 (N_T+R_TS_T)^{1-\vartheta_{\rm sp}}\\
&\qquad\times(C+DR_T)^{1-2\vartheta_{\rm sp}}
       +D^2N_TR_T\bigr]\|b\|_2^2.
\end{aligned}
\tag{3.15}
\]

Here $b^\flat(n'')$ is the dyadic-block slice of $\tilde b(n'')$
with $(n',a)=1$. The slice satisfies
$\|b^\flat(n'')\|_2\le\|\tilde b(n'')\|_2$.
Since $b$ is supported on the present dyadic block,
$b^\flat(n'')=\tilde b(n'')$ term by term.
Hence (3.15) is exactly (2.42)--(2.43).

The smooth derivatives in (3.14) obey the variable-based hypothesis
\eqref{p1:eq:2.39}, with the exponent $1-\epsilon_0$ displayed there:
each derivative is scaled by the corresponding positive argument
$c,d,n,r,s$ itself. In particular, the initial interval
$0<n\le N_T$ is not treated as a dyadic block of scale $N_T$.
Theorem 2.8 is used in the form of Section 2, (2.38)--(2.43).
It is proved in \S\ref{p1:sec:theorem-2-8-proof}.
Its spectral identity, including the continuous spectral measure, is exactly (2.14).
The I-2 estimate (Lemma~\ref{lem:typeI2}) also uses the composite-modulus Weil--Estermann bound \cite{Weil48,Estermann61} in \S\ref{sec:pi-nonII}.

The sum in (3.13) is a signed modulus sum; triangle inequalities over
source labels and boxes are applied outside it.

\subsection{Elementary preparation, moments, and SW}\label{sec:pi-sw}

\textbf{Divisor moments.} For integers $K,k\ge1$,

\[
\tau(n)^K\le\tau_{2^K}(n),\qquad
\sum_{n\le Y}\frac{\tau_k(n)}n\le(1+\log Y)^k.
\tag{3.16}
\]

At $p^v$, the first inequality is
$(v+1)^K\le\binom{v+2^K-1}{v}$.
The ratio of successive right/left quotients is at least 1.
This holds since $(1+t)^K\le1+(2^K-1)t$ for $0\le t\le1$, by convexity.
The second inequality follows by summing over the $k$ factors. Also
$\tau_j(n)\le\tau(n)^{j-1}$.

Suppose $Y^{1-1/k}\le H\le Y$. Select a largest factor in a
$k$-factorization. Then

\[
\sum_{Y<n\le Y+H}\tau_k(n)
 \le k\sum_{d\le(2Y)^{1-1/k}}\tau_{k-1}(d)(H/d+1)
 \le3kH(1+\log(2Y))^{k-1}.
\tag{3.17}
\]

Let $1\le u\le Y$ and $0<\varpi\le1$.
Split off prime factors larger than $Y^{\varpi/2}$.
Greedily group the rest into factors at most $Y^\varpi$.
There are at most $\lfloor2/\varpi\rfloor+1$ groups.
The large-prime part has at most $2/\varpi$ prime factors with multiplicity. Hence

\[
\tau(u)^K\le2^{2K/\varpi}
 \sum_{d\mid u,\ d\le Y^\varpi}
       \tau(d)^{K(\lfloor2/\varpi\rfloor+1)}.
\tag{3.18}
\]

For $u=|mn-a|\le X^3$, take $Y=X^3,\varpi=1/12$.
Put $h=(d,n)$. If $h\nmid a$, the congruence $d\mid mn-a$ has no solutions.
Otherwise it has at most $|J|h/d+1$ solutions on $J$.
Sum (3.18) and apply (3.16). For $|J|\le3M_I$ this gives

\[
\sum_{m\in J}\tau(|mn-a|)^K
 \le2^{24K}(1+\log X)^{2^{25K}}\{3M_I\tau(n)^{25K+1}+X^{1/4}\}.
\tag{3.19}
\]

Write $\widehat I_M$ for the integer support of $\psi_I$;
it is contained in an interval of length at most $3M_I$.
Sum (3.19) with $K=5$ over $n\in I_T$.
The $n$-sum of $\tau(n)^{126}\le\tau_{2^{126}}(n)$ is bounded by (3.17) and LEN.
Since $M_I\ge X^{1/2}$, the term $N_IX^{1/4}$ is absorbed by $M_IN_I$. Thus

\[
\sum_{n\in I_T}\sum_{m\in\widehat I_M}
                 \tau(|mn-a|)^5\ll M_IN_IL^{c_1},
\qquad c_1=2^{125}+2^{126}+2.
\tag{3.20}
\]
Here $c_1$ is an upper bound; the exponent $2^{125}+2^{126}-1$ already suffices.

For the prime source, Lemma~\ref{lem:aa37-typeII-lengths}
takes $T$ as the lower endpoint of the selected cell.
Under threshold (B.24), it proves (B.27), namely
$N_I\ge(2T)^{1-2^{-126}}$.
In (3.17) take $Y'=\lfloor T\rfloor$ and $H'=N_I$.
The integers in $I_T=(T,r_XT]$ of (B.22) are exactly
$(Y',Y'+N_I]$.
The stated bound implies $H'\ge(2Y')^{1-2^{-126}}$.
The mesh uses the lower endpoint of the dyadic shell, which is a different scale.
For the product source, LEN gives the lower length and the integer-interval convention.
Thus (3.17) applies to $\tau_{2^{126}}$, after increasing $N_0(\sigma)$ for each fixed
$\sigma>0$.
The same argument with $\varpi=\kappa/6$ treats
intervals of length at least $X^\kappa$.
It gives $\sum_m|\beta_m|^2\ll M_IL^c$ by elementary means; Shiu's estimate \cite{Shiu80} would also suffice.

\noindent\begin{minipage}{\linewidth}
\textbf{Finite Heath--Brown identity \cite{HB82}.} For $U\ge1$ and $n\le U^3$,

\[
\Lambda(n)=\sum_{j=1}^3(-1)^{j-1}\binom3j
 \sum_{n=m_1\cdots m_jn_1\cdots n_j,\ m_i\le U}
           \mu(m_1)\cdots\mu(m_j)\log n_1.
\tag{3.21}
\]
\end{minipage}

For the proof put $M(s)=\sum_{m\le U}\mu(m)m^{-s}$.
The series $1-\zeta M$ has no coefficients of index at most $U$. Hence
$-(\zeta'/\zeta)(1-\zeta M)^3$ has none of index at most $U^3$.
Expand the cube and use $-(\zeta'/\zeta)\zeta^j=(-\zeta')\zeta^{j-1}$.
Let $X$ be the dyadic shell of the tested coordinate, and $t\in(X,2X]$.
Take $U=2X^{1/3}$. Then (3.21) applies, since $t\le2X<U^3=8X$.
For the product source use $U_i=2X_i^{1/3}$ for the shell $X_i$ of each coordinate.
There are at most 24 variables.
Roughness factors exactly over the product. The identity

\[
1_{\mathbb P}(t)=\frac{\Lambda(t)}{\log t}
               -\sum_{k\ge2,\ t=p^k}\frac1k
\tag{3.22}
\]

\noindent has its proper-prime-power error bounded at source level in \S\ref{sec:pi-global}.

\textbf{Logarithms and boundaries.} On one short box replace each $1/\log p$
by its value at the box reference point. The pointwise error is
$O(L^{-C_{\rm box}-2})$ per coordinate. To sum the discrepancy
absolutely, use

\[
\sum_q|C_b(q)|\,|\mathcal S_t(q;a)|
 \ll_\sigma\tau(|t-a|)+L,\qquad t\ne a.
\tag{3.23}
\]

The first term counts divisors of $t-a$.
The second uses $\sum_{q\le X^{0.6}}1/\varphi(q)\ll L$.
The unrestricted labelled coefficient is $O_\sigma(1)$.
The number of prime labels is fixed,
and the small-prime part is disjoint from them.
For the prime source, the pointwise majorant in
Appendix~\ref{app:proof-aa37} gives (3.23) with
$F_a(n)=d(a-n)+2(1+L)$.
The source-wide logarithm-freezing estimate is (B.18):
$|E_{\rm LOG}|\ll_\sigma NL^{18-C_{\rm box}}$.
It is summed over tuples, each lying in one box, so its box-count
factor is 1. Here the logarithm is $L_N=\log N$, and its divisor-moment
input is (B.17). The box width is expressed in $L_X=\log X$.
Fixed powers of $L_X$ and $L_N$ are comparable.
For the product source, the corresponding aggregate bound is a
source condition stated in \S\ref{sec:pi-global}.

Write $L_N=\log N$. Consider the source cuts, the HB caps, and the monomial ordering and radial cuts of the product source.
On each dyadic source shell, each boundary region has relative width
$O(L_N^{-C_{\rm box}})$ in an active coordinate. Counting that
coordinate first and summing the geometric shell scales gives a total
of $O(NL_N^{5-C_{\rm box}}+N^{6/7}L_N^5)$ boundary indices, including
the fixed divisor multiplicity. An active coordinate can be chosen
among the positive-power factors, so the integer endpoint terms have a
fixed power saving. For the prime source, Cauchy--Schwarz with this
count and the source-wide second moment in (B.17) gives

\[
E_{\rm bdry}\ll N(\log N)^{c_R-C_{\rm box}/2}.
\tag{3.24}
\]

For the prime source in Theorem~\ref{thm:aa37-interior}, (3.24)
holds with $c_R=16$ by the global estimate (B.19) in
Appendix~\ref{app:proof-aa37}. This verifies (H6) for that source.

The prime-source estimate is invoked under
Theorem~\ref{thm:aa37-interior}, whose residue range supplies $a\ge N$.
This lower bound is compatible with the local upper bound:
$a\le N\log N/3\le(2/3)X(\log N)^{B+1}$ when
$X>N/(2(\log N)^B)$.
The prime-source cut and $a\ge N$ give
$|mn-a|\ge\varepsilon N/2$ on the doubled support.
Thus resonant boxes are empty and (3.10) holds.

\textbf{SW for the selected factor.} For every $A_{\rm SW}>0$, subinterval $I\subset I_T$,
integer $d\ge1$, and reduced class $l\bmod k$,

\[
\sum_{n\in I,\ n\equiv l(k),\ (n,d)=1}\alpha_n
 -\frac1{\varphi(k)}\sum_{n\in I,\ (n,dk)=1}\alpha_n
 \ll_{A_{\rm SW}} T\tau(d)L^{-A_{\rm SW}}.
\tag{3.25}
\]

Here and below the interval indicator is used only once.
Let $k>L^K$. After $K$ is chosen large, the trivial bound is
$O(T/k+1+T/\varphi(k))\ll TL^{-A_{\rm SW}}$.
The elementary bound $\varphi(k)\gg k/\log\log(k+2)$, or
$\varphi(k)\gg k^{1/2}$, suffices.
Let $k\le L^K$ and put $y_1=\exp(L^{2/3})$.
For sufficiently large $X$, $z_0>y_1$, so no selected integer
has largest prime below $y_1$.
Integers divisible by $p^2$, $p\ge y_1$, contribute $O(T/y_1)$.
Every other integer has a unique form $n=Pm$, with $P>P^+(m)$, $P\ge y_1$,
and $m\le2T/y_1$. For fixed $m$, its coefficient is
$\mu(m)\mu(P)=-\mu(m)$, $1$, or $\log(Pm)/L$.
The last has bounded variation on the allowed prime interval.
If $(m,dk)>1$, both relevant counts vanish,
or the coprimality with $d$ removes the term.
Otherwise $P\equiv l\bar m(k)$.
Apply prime SW and partial summation, uniformly at endpoints $P\ge y_1$.
The endpoint errors sum to $O_{A_{\rm SW}}(TL^{-A_{\rm SW}})$, since
$\log y_1=L^{2/3}$ and $\sum_{m\le2T/y_1}1/m\ll L$.
Excluding $P\mid dk$ costs at most
$O((\omega(d)+\omega(k))T/y_1)\ll_{A_{\rm SW}} T\tau(d)L^{-A_{\rm SW}}$.
Endpoint subtraction handles arbitrary $I$. This proves (3.25).
All errors are absolute, of size $O(TL^{-A_{\rm SW}})$.
Only the classical Siegel--Walfisz theorem is used for primes.
The proof is uniform for every $\zeta\ge z_0$.
A nonzero selected coefficient has all prime factors at least $\zeta$.
Raise the lower endpoint of the sum over the largest prime variable.
Its SW bound is unchanged. If $\zeta>2T$, the coefficient is zero.

\subsection{First dispersion with restoration before CRT}

Work first with the normalized convolution in (3.5).
Split all three factors dyadically.
Apply the triangle inequality before squaring.
There are $O(L^3)$ triples, including the outer factor. For one triple put

\[
\lambda_r=\sum_{r_3r_2=r}\gamma_3(r_3)\gamma_2(r_2),
\quad r_3\asymp R,\ r_2\asymp S.
\tag{3.26}
\]

Since $r$ is $u_0$-rough and at most $X^{0.6}$,
$|\lambda_r|\le\tau(r)=O_\sigma(1)$. The arbitrary small-prime factor
of C1 remains in the outer $q$, as the only smooth divisor.
Let $\psi_0$ be a smooth nonnegative bump, equal to 1 on $[1,2]$,
supported on $[1/2,5/2]$. Cauchy--Schwarz in $(q,m)$ gives

\[
\begin{aligned}
&|S|^2\ll QM_IL^c S',\quad\\
&S'=\sum_q\psi_0(q/Q)1_{(q,a)=1}\sum_m\psi_I(m)
 \left|\sum_r\lambda_r\sum_n\alpha_n\mathcal S_{mn}(qr;a)\right|^2.
\end{aligned}
\tag{3.27}
\]

Write $t=(q,r_1,r_2,n_1,n_2)$, $k_i=mn_i-a$, and

\[
c(t)=\psi_0(q/Q)1_{(q,a)=1}\lambda_{r_1}\bar\lambda_{r_2}
                       \alpha_{n_1}\bar\alpha_{n_2},
\quad \omega(t)=1_{(n_1,qr_1)=(n_2,qr_2)=1}.
\tag{3.28}
\]

Expanding the square gives $S'=S_1-2\Re S_2+S_3$, where

\[
\begin{aligned}
S_3&=\sum_t c(t)\sum_m\psi_I(m)1_{qr_1\mid k_1}1_{qr_2\mid k_2},\\
S_2&=\sum_t c(t)\sum_m\psi_I(m)
       1_{qr_1\mid k_1}1_{(mn_2,qr_2)=1}/\varphi(qr_2),\\
S_1&=\sum_t c(t)\sum_m\psi_I(m)
       1_{(mn_1,qr_1)=1}1_{(mn_2,qr_2)=1}/(\varphi(qr_1)\varphi(qr_2)).
\end{aligned}
\tag{3.29}
\]

Let $\mathcal T$ be the tuple set. Let $\mathfrak B$ be its union
of cross-copy gcd failures

\[
(n_1,n_2)>1,\ (r_1,r_2)>1,\ (n_1,r_2)>1,
\ (n_2,r_1)>1,\ (q,r_1r_2)>1,
\tag{3.30}
\]

and let $\mathfrak D=\{|n_1-n_2|<TL^{-C_M}\}$.
Set $\mathfrak G'=\mathcal T\setminus(\mathfrak B\cup\mathfrak D)$.
On every tuple subset $V$ define the formal sums

\[
\begin{aligned}
MT(V)&=\sum_{t\in V}c(t)\omega(t)
                \frac{M_I\widehat\phi_0(0)}{qr_1r_2\varphi(q)},\\
\mathcal Z(V)&=\sum_{t\in V}c(t)\omega(t)1_{n_1\equiv n_2(q)}
                \frac{M_I\widehat\phi_0(0)}{qr_1r_2}.
\end{aligned}
\tag{3.31}
\]

\textbf{Restoration lemma.} Choose $A_E,C_\Sigma,C_M$ in the order of \S\ref{sec:pi-global}. Then

\[
S'=M_IE_2^*+\mathcal Z(\mathcal T)-MT(\mathcal T)+O(E_{\rm GR}),
\tag{3.32}
\]
\[
\begin{aligned}
E_{\rm GR}\ll{}&L^{3C_\Sigma+c}M_IN_I\left(\frac{N_I}{z_0Q}+1\right)
 +C_\epsilon X^\epsilon M_IN_I\left(\frac{N_I}{u_0Q}+1\right)\\
&+\frac{M_IN_I^2}{QL^{A_E}}.
\end{aligned}
\tag{3.33}
\]

We first prove the positive estimates needed for CRT.
For $S_i$, let $\Upsilon_i(V)$ be the positive majorant.
It replaces $c(t)$ by $|c(t)|$ and deletes coprimalities.
Retain the two congruences for $\Upsilon_3$, one for $\Upsilon_2$.
Let $|MT|(V),|\mathcal Z|(V)$ be the positive versions of (3.31).
We write $\Upsilon_{\rm exc}$ for the positive contribution of
$(S_3)$ from the complement of $\Sigma^*$ defined in (3.37).

The total-variation progression estimate is

\[
\left|\sum_{m\equiv b(k)}f(m)-k^{-1}\int f\right|\le V(f).
\tag{3.34}
\]

Compare $f(b+jk)$ to its cell average on
$[b+jk,b+(j+1)k]$. Bound the difference by the cell variation,
and sum. Möbius expansion of the coprimality then gives
$\sum_{(m,k)=1}\psi_I=(\varphi(k)/k)M_I\widehat\phi_0(0)
+O(\tau(k))$.

From $n/\varphi(n)=\sum_{d\mid n}\mu^2(d)/\varphi(d)$ we get

\[
\sum_{r\le Y}\frac1{\varphi(r)}\ll\log(2Y),\quad
\sum_{r\le Y,p\mid r}\frac1{\varphi(r)}\ll\frac{\log(2Y)}{p-1},
\]
\[
\sum_{q\asymp Q}\frac1{\varphi(q)^2}\ll L^2/Q,\qquad
\sum_{q\asymp Q,p\mid q}\frac1{\varphi(q)^2}\ll L^2/(pQ).
\tag{3.35}
\]

For $\Upsilon_1,\Upsilon_2,|MT|$ the weight is bounded by
$O_\sigma(M_I/(\varphi(q)^2\varphi(r_1)\varphi(r_2)))$, since
$qr_1\ll QRS<M_I$. For $|\mathcal Z|$ it is
$O_\sigma(M_I/(qr_1r_2))1_{n_1\equiv n_2(q)}\omega(t)$.
These two shapes prove

\[
\Upsilon_1(\mathfrak B)+\Upsilon_2(\mathfrak B)+|MT|(\mathfrak B)+|\mathcal Z|(\mathfrak B)
 \ll L^cM_IN_I\left(\frac{N_I}{z_0Q}+1\right).
\tag{3.36}
\]

For the first failure of (3.30), fix $n_1$.
Choose one of its at most $L$ primes $p\ge z_0$.
Count $n_2$ in multiples of $p$. The first shape gives $N_I/z_0+1$.
For the weight of $|\mathcal Z|$, $(n_1,q)=1$.
So the extra congruence gives a class modulo $pq$,
with $N_I/(pq)+1$ choices.
For a common prime $p\ge u_0$ of $r_1,r_2$, fix $r_1$.
Sum the $p$-divisible $r_2$ weight using (3.35).
For $p\mid(n_1,r_2)$, fix $n_1$ and do the same weighted $r_2$ sum.
For $p\mid(n_2,r_1)$, fix $r_1$.
Then count $n_2$ in multiples of $p$, or modulo $pq$ for the weight of $|\mathcal Z|$.
For $p\mid(q,r_i)$, use the last estimate of (3.35).
This covers all five failures, and $u_0\ge z_0$ gives (3.36).

Split the unrestricted $(t,m)$-sum at

\[
\Sigma^*=\{\tau(|k_1|)\tau(|k_2|)<L^{C_\Sigma}\}.
\tag{3.37}
\]

The complement is estimated on all tuples, coprime or not.
Put $u=\tau(|k_1|)$ and $v=\tau(|k_2|)$.
The exceptional indicator is at most $uvL^{-C_\Sigma}$.
The two divisor choices $r_i\mid k_i$ cost at most $uv$.
There are also at most $u$ choices for $q\mid k_1$.
The resulting weight is $u^3v^2$, of degree five.
To keep the progression restriction while counting, split into $u\ge v$
and $v>u$.
On the first part $u^3v^2\le u^5$.
For each fixed $(m,n_1)$, sum over the at most $u$ divisors $q$ first.
The condition $q\mid k_1$ and $(q,a)=1$ implies $(m,q)=1$.
Thus $q\mid k_2$ forces $n_2\equiv n_1\pmod q$.
There are at most $4N_I/Q$ such $n_2$, by (3.6).
This bounds the first part by $L^{-C_\Sigma}(4N_I/Q)\sum u^5$.
Exchanging the two copies $(n_1,r_1)\leftrightarrow(n_2,r_2)$ gives the same bound for the second part.
Consequently

\[
\Upsilon_{\rm exc}\ll L^{-C_\Sigma}\frac{N_I}{Q}
 \sum_{n\in I_T,m\in\widehat I_M}\tau(|mn-a|)^5
 \ll \frac{M_IN_I^2}{Q}L^{c_1-C_\Sigma}.
\tag{3.38}
\]
The factor $4$ from $4N_I/Q$ and the factor $2$ from the two parts are
absorbed in the implied constant.

For a modulus-related failure, fix $(m,n_1)$. The number of
$(q,r_1)$ with $qr_1\mid k_1\ne0$ is
$\tau_3(|k_1|)\le\tau(|k_1|)^2\ll_\epsilon X^{\epsilon/3}$.
A common prime $p\ge u_0$ comes from $r_1$, $n_1$, or $q$.
It forces $n_2$ into one class modulo $pq$.
For $p\mid r_1,r_2$ use $p\nmid m$.
For $p\mid n_1,r_2$ use $mn_2\equiv a(p)$.
For $p\mid n_2,r_1$ use $n_2\equiv0(p)$.
If $p\mid q,r_2$, then $pq\mid qr_2\mid k_2$ gives the class directly.
If $p\mid q,r_1$ only, exchange the copies.
There are $O(L)$ candidate primes and $O(N_I/(u_0Q)+1)$ choices in each class.
There are at most $\tau(|k_2|)$ remaining choices of $r_2$. Hence

\[
\Upsilon_3(\mathfrak B_{\rm modulus})
 \ll_\epsilon X^\epsilon M_IN_I(N_I/(u_0Q)+1).
\tag{3.39}
\]

For $(n_1,n_2)>1$ on $\Sigma^*$, there are at most $L^{2C_\Sigma}$
pairs $(q,r_1)$. There are at most $L$ primes $p\mid n_1$, $p\ge z_0$.
There are at most $2N_I/(z_0Q)+1$ choices for $n_2$.
There are at most $L^{C_\Sigma}$ choices for $r_2$. Thus

\[
\Upsilon_3(\mathfrak B_n\cap\Sigma^*)
 \ll L^{3C_\Sigma+1}M_IN_I(N_I/(z_0Q)+1).
\tag{3.40}
\]

On the near-diagonal the same count replaces $N_I/z_0$ by
$TL^{-C_M}$. The density shapes and (3.35) give the same or smaller
bound. Therefore

\[
\begin{aligned}
&\Upsilon_3(\mathfrak D\cap\Sigma^*)+\Upsilon_1(\mathfrak D)+\Upsilon_2(\mathfrak D)
 +|MT|(\mathfrak D)+|\mathcal Z|(\mathfrak D)\\
&\qquad\ll L^{3C_\Sigma+c}M_IN_I(TL^{-C_M}/Q+1).
\end{aligned}
\tag{3.41}
\]

Use the common constants from \S\ref{sec:pi-global}:
$C_\Sigma=A_E+2c_*+3$ and
$C_M=5C_\Sigma+A_E+C_{\rm box}+c_*+3$.
They satisfy $C_\Sigma>A_E+c_1+2$ and
$C_M>3C_\Sigma+C_{\rm box}+A_E+c$.
To check this, divide by the last term of (3.33). After this division, the right side
of (3.41) is at most
\[
L^{3C_\Sigma+c+A_E}
\left(\frac{T L^{-C_M}}{N_I}+\frac{Q}{N_I}\right)
\ll L^{3C_\Sigma+c+A_E+C_{\rm box}-C_M}
   +L^{3C_\Sigma+c+A_E}X^{-\sigma/3}=o(1).
\]
Here LEN supplies $T/N_I\ll L^{C_{\rm box}}$, while (3.6)
supplies $Q/N_I\le X^{-\sigma/3}$. The chosen $C_M$ makes
the first logarithmic exponent negative; the second term tends to zero
for fixed $\sigma$. The same target absorbs (3.38), since
$C_\Sigma>A_E+c_1+2$.
This removes bad tuples and the near-diagonal from \textbf{all} of
$S_1,S_2,S_3,Z,MT$ before their main terms are compared.

Now use CRT on $\mathfrak G'$. Let
$Q_1=qr_1,Q_2=qr_2,g=(r_1,r_2)$.
The two congruences have one class modulo $[Q_1,Q_2]=qr_1r_2/g$ precisely when
$(n_1,Q_1)=(n_2,Q_2)=1$ and $n_1\equiv n_2(qg)$.
To see this, equate $a\bar n_1$ and $a\bar n_2$ modulo $(Q_1,Q_2)=qg$.
On $\mathfrak G'$, $g=1$ and $q,r_1,r_2$ are pairwise coprime.

\textbf{Fixed cutoff.} On the smoothed support put

\[
k(t)=qr_1r_2\le K_*=40QR^2S^2,\quad
H_2=QR^2S^2X^\eta/M_I,\quad H_*=40H_2.
\tag{3.42}
\]

These are fixed for the entire shell triple. Poisson gives exactly

\[
\begin{aligned}
E_2^*={}&\sum_{t\in\mathfrak G'}
 \frac{c(t)\omega(t)1_{n_1\equiv n_2(q)}}{k(t)}\\
&\times\sum_{1\le|h|\le H_*}\widehat\phi_0(hM_I/k(t))
                      e(h(b(t)-m_0)/k(t)),\\
S_3(\mathfrak G')={}&\mathcal Z(\mathfrak G')+M_IE_2^*+O(X^{-90}).
\end{aligned}
\tag{3.43}
\]

Put $\lambda=M_I/k(t)$. This is the Poisson prefactor of one
progression sum. In (3.43) it appears as the factor $M_I$ in front of
$E_2^*$ and the factor $1/k(t)$ inside $E_2^*$.
Integration by parts gives
$|\widehat\phi_0(u)|\le C_j|u|^{-j}$.

Here one may take $C_j=(2\pi)^{-j}\|\phi_0^{(j)}\|_1$.
It is independent of the box and the tuple.
The integer $j=\lceil104/\eta\rceil+1$ is fixed once $\sigma$ is fixed.
Choose $X_0$ after these seminorms and all fixed coefficient parameters.
So the tail bound is uniform over the whole box family.
If $H_*<1$, then $\lambda>X^\eta$, since $\lambda H_*\ge X^\eta$.

The tail of one progression
sum is at most $C'_j(\lambda H_*)^{1-j}$ if $H_*\ge1$, and at
most $C'_j\lambda^{1-j}$ otherwise. In either case it is
$O(X^{-\eta(j-1)})$, since $\lambda H_*\ge X^\eta$.
Take $j=\lceil104/\eta\rceil+1$. After increasing $X_0$, the
weighted tail per tuple is at most $X^{-103}$. There are at most
$X^5$ tuples. This proves (3.43).
The cap (3.42) is the same for all tuples.
So it survives the Mellin separation of the inner variable below.

On $\mathfrak G'$, (3.34) gives

\[
S_1(\mathfrak G')=MT(\mathfrak G')+O(N_I^2L^c/Q).
\tag{3.44}
\]

For $S_2$, retain $(m,r_2)=1$. Möbius expansion over $e\mid r_2$
gives one class modulo $qr_1e$, so

\[
\sum_m\psi_I(m)1_{mn_1\equiv a(qr_1)}1_{(m,r_2)=1}
 =\frac{M_I\widehat\phi_0(0)\varphi(r_2)}{qr_1r_2}
   +O(\tau(r_2)).
\tag{3.45}
\]

Divide by $\varphi(qr_2)=\varphi(q)\varphi(r_2)$.
The main term is again $MT$.
The summed error is $O(RSN_I^2L^c)$.
It is absorbed using $QRS<M_IX^{-89\sigma}$.
Exchanging the two copies shows that $MT(\mathfrak G')$ is real.
Complete $\mathcal Z$ and $MT$ formally to $\mathcal T$, with errors bounded by (3.36), (3.41).
This proves (3.32).
CRT and the zero-mode identity are used on $\mathfrak G'$ only.

\subsection{Main-term cancellation and the full GR error}

Put $\alpha^{(r)}_n=\alpha_n1_{(n,r)=1}$ and

\[
\Delta(\xi;q,b)=\sum_{n\equiv b(q)}\xi_n-
                  \varphi(q)^{-1}\sum_{(n,q)=1}\xi_n.
\tag{3.46}
\]

Expanding the product of these two centered counts gives exactly

\[
\begin{aligned}
\mathcal Z(\mathcal T)-MT(\mathcal T)
={}&\sum_{q,r_1,r_2}\frac{\psi_0(q/Q)1_{(q,a)=1}
               \lambda_{r_1}\bar\lambda_{r_2}M_I\widehat\phi_0(0)}{qr_1r_2}\\
&\qquad\sum_{b\bmod q}^{*}\Delta(\alpha^{(r_1)};q,b)
                  \overline{\Delta(\alpha^{(r_2)};q,b)}.
\end{aligned}
\tag{3.47}
\]

For the norm-normalized condition \cite[(A2)]{BFI86}, apply (3.25) to $dr$, with
$\tau(dr)\le\tau(d)O_\sigma(1)$. If
$\|\alpha^{(r)}\|_2^2\le N_IL^{-A_3-1}$, the elementary estimate

\[
\sum_{q\le3Q}\sum_b^*|\Delta(\alpha^{(r)};q,b)|^2
 \ll (T\log(3Q)+3Q)\|\alpha^{(r)}\|_2^2
 \ll N_ITL^{-A_3}
\tag{3.48}
\]

\noindent suffices, since $\log(3Q)\le L$ is absorbed by the factor $L^{-1}$.
If $\|\alpha^{(r)}\|_2^2>N_IL^{-A_3-1}$, then $\|\alpha^{(r)}\|_2(2T)^{1/2}\gg TL^{-(C_{\rm box}+A_3+1)/2}$.
Then (3.25) with exponent
$A_{\rm BFI}+(C_{\rm box}+A_3+1)/2$ implies \cite[(A2)]{BFI86}.
The condition $Q<TX^{-29\sigma}$ implies the condition
$3Q\le T(\log T)^{-B_1}$ of \cite[Theorem~0(a)]{BFI86} eventually,
which then gives (3.48) in this case too.
Apply $|uv|\le(|u|^2+|v|^2)/2$ in (3.47), and sum
the reciprocal $r_i$ weights. This gives

\[
|\mathcal Z(\mathcal T)-MT(\mathcal T)|
 \ll M_IN_I^2L^{C_{\rm box}+2-A_3}/Q.
\tag{3.49}
\]

Take $A_3=A_E+C_{\rm box}+2$.
Increase it by a fixed constant to absorb any endpoint logarithm.
Combine (3.27), (3.32), (3.33), (3.49), with $\epsilon=\ell_0\sigma^2/8$.
For a requested local precision $A'>0$, set $A_{\rm work}=A'+4$
before choosing $A_E,C_\Sigma,C_M$ as in \S\ref{sec:pi-global}.
The resulting restoration contributions to $|S|$ are at most

\[
XL^{-A_{\rm work}}+XL^Cz_0^{-1/2}
 +XL^C(Q/N_I)^{1/2}
 +XX^{\epsilon/2}L^C\{u_0^{-1/2}+(Q/N_I)^{1/2}\}.
\tag{3.50}
\]

Both terms in braces are kept.
Their power savings are at least $7\ell_0\sigma^2/16$ and
$\sigma/6-\ell_0\sigma^2/16$, respectively.
Each exceeds $c_\sigma$, and $z_0^{-1/2}$ beats every logarithm.
More explicitly, for fixed $A_{\rm work}$ and sufficiently large $X$,
$XL^Cz_0^{-1/2}\ll XL^{-A_{\rm work}}$, because
$z_0^{-1/2}=X^{-1/(2(\log\log X)^3)}$ beats every fixed logarithm.
Thus this term joins the first term in (3.51), rather than the
fixed-power-saving term.
The separate term $XL^C(Q/N_I)^{1/2}$ in (3.50) has saving
$\sigma/6>c_\sigma$, by $N_I\ge QX^{\sigma/3}$ in (3.6),
and is therefore absorbed in the displayed power-saving term.
Increase $X_0$ after the fixed exponents. This gives

\[
|S|\ll XL^{-A_{\rm work}}
      +M_IQ^{1/2}L^c|E_2^*|^{1/2}+X^{1-c_\sigma}L^C.
\tag{3.51}
\]

This completes the first dispersion.
A later saving for $E_2^*$ changes only the middle term.

\subsection{Reciprocity, the complete mask, and the second square}

Split the combined products again as $r_i\leftarrow r_is_i$, where
$r_i\asymp R,s_i\asymp S$. The remaining support is

\[
(q,r_1s_1r_2s_2)=(r_1s_1,r_2s_2)=(n_1,n_2)=1,
\]
\[
(n_1,r_2s_2)=(n_2,r_1s_1)=1,
\quad(n_i,qr_is_i)=1,
\quad |n_1-n_2|\ge TL^{-C_M}.
\tag{3.52}
\]

Solving the two CRT congruences and using
$\bar u/v+\bar v/u\equiv1/(uv)\pmod1$ for $(u,v)=1$ gives

\[
\frac{ah\overline{n_1r_2s_2}}{qr_1s_1}
 +\frac{ah\overline{n_2qr_1s_1}}{r_2s_2}
 \equiv\frac{ah(n_1-n_2)\overline{n_2qr_1s_1}}{n_1r_2s_2}
       +\frac{ah}{qr_1s_1r_2s_2n_1}\pmod1.
\tag{3.53}
\]

All inverses exist by (3.52). Put $f=(n_1-n_2)/q\ne0$.
Then the outer tuple $\omega=(f,n_1,n_2,r_1,r_2,s_2)$ determines $q=(n_1-n_2)/f$.
The modular phase becomes

\[
e\left(\frac{ahf\overline{n_2r_1s_1}}{n_1r_2s_2}\right).
\tag{3.54}
\]

Split $f\ne0$ by sign and dyadic intervals
$f\in\varsigma_f(F,2F]$, starting at $F=1/2$. Split $h\ne0$ by
sign and intervals $H_0<|h|\le\min(2H_0,H_*)$, starting at
$H_0=1/2$. These disjoint intervals include $f,h=\pm1$ exactly
once. Write $\varsigma_f,\varsigma_h\in\{\pm1\}$ for the signs of $f$ and $h$.
The cap on the last, shortened $h$-block depends on $h$ alone.
Always $F\ll T/Q$.

The nonmodular factor is exactly

\[
\widehat\phi_0\left(\frac{hM_I}{qr_1s_1r_2s_2}\right)
 e\left[\frac{h}{qr_1s_1r_2s_2}
                   \left(\frac a{n_1}-m_0\right)\right].
\tag{3.55}
\]

Its dependence on $(h,s_1)$ is only through $|h|/s_1$. For
$v=|h|S/(H_0s_1)\in[1/2,2]$, choose a fixed smooth cutoff.
It equals 1 there and is supported in $(1/4,4)$. Rewrite (3.55):

\[
\begin{gathered}
W_{Y(\omega)}(v)=w(v)\widehat\phi_0(Y_2v)e(Y_1v),\\
Y_1=\frac{\varsigma_hH_0(a/n_1-m_0)}{Sqr_1r_2s_2},\qquad
Y_2=\frac{\varsigma_hH_0M_I}{Sqr_1r_2s_2}.
\end{gathered}
\tag{3.56}
\]

Since $H_0\le H_*=40H_2$,

\[
|Y_1|+|Y_2|\ll (1+a/X)L^{C_{\rm box}}X^\eta
 \ll L^{C_{\rm box}+K_{\rm scale}}X^\eta=:L_{\rm osc}.
\tag{3.57}
\]

Mellin inversion separates the ratio. Put
$m(t)=\sup_\omega|\widetilde W_{Y(\omega)}(it)|$.
The zeroth and second logarithmic derivative estimates give

\[
m(t)\ll\min\{1,L_{\rm osc}^2/(1+|t|)^2\},\qquad
\int_\mathbb Rm(t)\,dt\ll L_{\rm osc}.
\tag{3.58}
\]

If $m(t)>0$, divide the outer coefficient by $m(t)$.
The norm bounds below are for this normalized coefficient.
The factor $m(t)$ is restored before the $t$-integration.
If $m(t)=0$, the integrand is zero.
So the integrable decay of $m$ is kept.

The full inner coprimality mask is

\[
(s_1,ac)=1,\qquad(s_1,J_\omega)=1,
\quad c=n_1r_2s_2,\quad J_\omega=qn_2.
\tag{3.59}
\]

The second condition combines $(s_1,q)=1$ and $(s_1,n_2)=1$.
To separate the second, split $\mathcal P_X=[u_0,X^3]\cap\mathbb P$ at random.
Each prime goes independently and uniformly into one of two classes $B_\chi,R_\chi$.
For an integer $m$, let
$\mathcal P_X(m)=\{p\in\mathcal P_X:p\mid m\}$ and
$\nu_X(m)=|\mathcal P_X(m)|$. If all prime factors of $s$ lie in
$\mathcal P_X$, then exactly

\[
1_{(s,J)=1}=\mathbb E_\chi
 [2^{\nu_X(s)+\nu_X(J)}
  1_{\mathcal P_X(s)\subseteq B_\chi}
  1_{\mathcal P_X(J)\subseteq R_\chi}].
\tag{3.60}
\]

A common prime makes the event impossible.
Otherwise the event has probability $2^{-\nu_X(s)-\nu_X(J)}$.
Every prime of a supported $s_1$ lies in $\mathcal P_X$.
So $(s_1,J_\omega)=1$ depends only on the primes of $J_\omega$ in $\mathcal P_X$.
Since $s_1\le X^{0.6}$ and $J_\omega\le X^{0.94}$,
one may take

\[
K_s=\lceil0.6/(\ell_0\sigma^2)\rceil,\qquad
K_J=\lceil1/(\ell_0\sigma^2)\rceil.
\tag{3.61}
\]

Replace the inner coefficient by
$\lambda_s^\chi=\lambda_s2^{\nu_X(s)-K_s}
1_{\mathcal P_X(s)\subseteq B_\chi}$
and each outer representation coefficient by
$A_\omega^\chi=A_\omega2^{\nu_X(J_\omega)-K_J}
1_{\mathcal P_X(J_\omega)\subseteq R_\chi}$.
Both multipliers have modulus at most 1.
The original sum is $2^{K_s+K_J}\mathbb E_\chi W_\chi$.
This factor is $O_\sigma(1)$, and $\mathbb E_\chi$ is an average.
Fix $\chi$ before aggregating representations.
The partition acts on $\lambda_s$ and $A_\omega$ only, so (3.25) for $\alpha_n$ holds for each $\chi$.

Set $b=n_2r_1,c=n_1r_2s_2$, and

\[
I(b,c,f)=\sum_{s\asymp S,(s,ac)=1}\sum_h
             \lambda'_{s,t}\eta_{h,t}
             e(ahf\overline{bs}/c),
\tag{3.62}
\]

where $\lambda'_{s,t}=(S/s)\lambda_s^\chi(s/S)^{it}$, and
$\eta_{h,t}=(|h|/H_0)^{-it}1_{|h|\le H_*}$ on the chosen signed
block. Both are 1-bounded. Dyadically localize $b\asymp D,c\asymp C$.
The outer support has $(b,c)=(b,p_0)=1$. So the nonnegative majorant

\[
\rho=\psi_0(b/D)\psi_0(c/C)1_{(b,c)=(b,p_0)=1}
                                   1_{f\in\varsigma_f(F,2F]}
\tag{3.63}
\]

\noindent equals 1 on it. Cauchy--Schwarz now gives

\[
|W_{\chi,t}|\le\left(\sum_{b,c,f}|A_{b,c,f}^\chi|^2\right)^{1/2}
               W_4^{1/2},\quad W_4=\sum_{b,c,f}\rho|I|^2.
\tag{3.64}
\]

Fix $(b,c,f)$. Factor $b=n_2r_1,c=n_1r_2s_2$; then $q$ is determined.
So there are at most $\tau(b)\tau_3(c)$ outer representations.
For fixed $n_1,n_2$ the permitted $q\mid n_1-n_2$ have divisor multiplicity.
So there are $O(T^2R^2S)X^{o(1)}$ representations in total.
A normalized representation coefficient
has modulus $O_\sigma((QR^2S^2)^{-1})$. Summing over representations with Cauchy--Schwarz,

\[
\left(\sum|A^\chi|^2\right)^{1/2}
 \ll X^{o(1)}\frac{T}{QRS^{3/2}}.
\tag{3.65}
\]

This bound holds for the coefficients $A^\chi$ of each class $\chi$.

Expanding $|I|^2$, put

\[
z=ss',\quad k=h_1s'-h_2s,\quad y=fk.
\tag{3.66}
\]

The terms with $k=0$ have at most $X^{o(1)}H_*S$ choices of
$(s,s',h_1,h_2)$. To see this, write $(s,s')=g(u,v)$, $(u,v)=1$, and
$(h_1,h_2)=j(u,v)$. Sum $(g,j)$ over their ranges. Thus

\[
|W_{k=0}|\ll X^{o(1)}H_*S DFC
 \ll (T^2R^2S^3)X^{-8\sigma+\eta+o(1)},
\tag{3.67}
\]

\noindent where the final inequality is $T^2R^2S<X^{1-8\sigma}$.
For $k\ne0$, split $y=\varepsilon_y m$, $m>0$, and $z$ into
$Z\in\{S^2,2S^2\}$. The coefficient independent of $(b,c)$ is

\[
\begin{aligned}
\beta^{\chi,\varepsilon_y}_{z,m}={}&1_{(z,ap_0)=1}
 \sum_{ss'=z}\lambda'_{s,t}\bar\lambda'_{s',t}
 \sum_{\substack{f\in\varsigma_f(F,2F],\ h_i\ {\rm in\ block}\\
                  f(h_1s'-h_2s)=\varepsilon_y m}}
                     \eta_{h_1,t}\bar\eta_{h_2,t}.
\end{aligned}
\tag{3.68}
\]

Put $Y:=H_*SF$. By (3.70), $|y|\le8FH_*S$.
So the coefficient is supported on $1\le m\le8Y$.
The off-diagonal sum is exactly

\[
\sum_{z,m}\beta_{z,m}^{\chi,\varepsilon_y}
 \sum_{b,c:(b,c)=(b,p_0)=(z,c)=1}
      \psi_0(b/D)\psi_0(c/C)e(\varepsilon_y a m\overline{zb}/c).
\tag{3.69}
\]

\subsection{The W-prime norm, with endpoints}

For $S\ge1$, $1\le H\le S$, $F\ge1/2$, and $a\ne0$, let

\[
\begin{aligned}
B_{z,y}=\#\{&s_i\in(S,2S],\ (s_1s_2,a)=1,\\
&1\le|h_i|\le H,\ f\in(F,2F]:\\
&s_1s_2=z,\ f(h_1s_2-h_2s_1)=y\ne0\}.
\end{aligned}
\tag{3.70}
\]

For every nonzero integer $v$, define its $a$-supported part by
$v_a=\prod_{p\mid a}p^{v_p(|v|)}$.
In particular $y_a=\prod_{p\mid a}p^{v_p(|y|)}$.
Set

\[
\begin{gathered}
\mathcal T_0=\max\tau(z)\tau(|y|),\qquad
P_0=\prod_{p\mid a}(1-p^{\vartheta_{\rm sp}-1})^{-1},\\
P_1=P_0^2,\qquad P_2=\prod_{p\mid a}(1-p^{\vartheta_{\rm sp}-1/2})^{-1},
\end{gathered}
\tag{3.71}
\]

where the maximum is over $z\le4S^2,0<|y|\le8FHS$.
Then

\[
\sum_{z,y\ne0}B_{z,y}^2(y_a)^{\vartheta_{\rm sp}}
 \le\mathcal T_0\,2FP_0H^2S^2
             [16P_1+64(8+\log(2S))P_0+32P_2].
\tag{3.72}
\]

\emph{Proof.} For fixed $s=(s_1,s_2)$, let
$N_s(k)=\#\{h:1\le|h_i|\le H,h_1s_2-h_2s_1=k\}$.
Each $(z,y)$ has at most $\tau(z)\tau(|y|)$ choices of $(s,f)$.
Therefore Cauchy--Schwarz and $(fk)_a=f_ak_a$ reduce the left side to

\[
\mathcal T_0\left(\sum_{f\le2F}f_a^{\vartheta_{\rm sp}}\right)Q_0,
\qquad Q_0=\sum_s\sum_{k\ne0}k_a^{\vartheta_{\rm sp}} N_s(k)^2,
\quad \sum_{f\le2F}f_a^{\vartheta_{\rm sp}}\le2FP_0.
\tag{3.73}
\]

For the last inequality, majorize $f_a^{\vartheta_{\rm sp}}$ by
$\sum_{d\mid a^\infty,d\mid f}d^{\vartheta_{\rm sp}}$.
Then sum $\lfloor2F/d\rfloor\le2F/d$. It holds for every $F>0$.

Put $g=(s_1,s_2)$, $s=g(u,v)$, $S'=S/g$.
A fibre of $h_1s_2-h_2s_1$ has at most $1+2H/S'$ points.
It has at most one when $S'\ge2H$. Expand

\[
k_a^{\vartheta_{\rm sp}}=\sum_{d\mid a^\infty,d\mid k}\varphi_{\vartheta_{\rm sp}}(d),
\quad\varphi_{\vartheta_{\rm sp}}(p^j)=p^{j\vartheta_{\rm sp}}(1-p^{-\vartheta_{\rm sp}}),
\quad0\le\varphi_{\vartheta_{\rm sp}}(d)\le d^{\vartheta_{\rm sp}}.
\tag{3.74}
\]
Here $\varphi_{\vartheta_{\rm sp}}$ is multiplicative with
$\varphi_{\vartheta_{\rm sp}}(1)=1$; the prime-power formula is for $j\ge1$.

If $A_d(s)$ counts $h\in[-H,H]^2$ with $k\ne0,d\mid k$, then

\[
Q_0\le\sum_{d\mid a^\infty,d\le4HS}d^{\vartheta_{\rm sp}}
 \left[\sum_s A_d(s)+\sum_{s:S'<2H}\frac{2H}{S'}A_d(s)\right].
\tag{3.75}
\]

For the first sum write $h=e w$, $w$ primitive,
$\|w\|_\infty=j\le H$. There are at most $8j$ choices of $w$,
and $e\le H/j$. The congruence is
$d_e\mid w_1s_2-w_2s_1$, $d_e=d/(d,e)$.
Since $(s_1s_2,d_e)=1$, a contributing primitive vector has
$(w_1,d_e)=1$. The case $w_1=0$ forces $d_e=1$. There are at
most $2S$ choices for $s_1$, then $S/d_e+1$ for $s_2$.
By $\sum_{e\le E}(d,e)\le\tau(d)E$, this proves

\[
\sum_sA_d(s)\le16\tau(d)H^2S^2/d+16H^2S.
\tag{3.76}
\]

For the second sum $(d,g)=1$.
At most $8HS'/d$ nonzero levels of $h_1v-h_2u$ are divisible by $d$.
Each has at most $4H/S'$ points. Thus $A_d(s)\le32H^2/d$.
For each $g$, sum the at most $(S/g+1)^2$ pairs.
Use the exact inequality

\[
\sum_{g\le2S}\frac{(S/g+1)^2}{S/g}
 =\sum_{g\le2S}(S/g+2+g/S)\le S(8+\log(2S)).
\tag{3.77}
\]

So before the $d$-sum, the second term of (3.75) is at most
$64(8+\log(2S))H^3S/d\le64(8+\log(2S))H^2S^2/d$, using $H\le S$. Finally

\[
\begin{aligned}
\sum_{d\mid a^\infty}\tau(d)d^{\vartheta_{\rm sp}-1}&=P_1,\\
\sum_{d\mid a^\infty}d^{\vartheta_{\rm sp}-1}&=P_0,\\
\sum_{d\le4HS,d\mid a^\infty}d^{\vartheta_{\rm sp}}&\le(4HS)^{1/2}P_2.
\end{aligned}
\tag{3.78}
\]

The last contribution of (3.76) is at most $32H^2S^2P_2$, since
$H\le S$. This proves (3.72), also for real nonintegral $(S,H)$.

Uniformly for $|a|\le X^2$, elementary prime ordering gives

\[
\begin{gathered}
P_0,P_1,P_2\le\exp(C(\log X)^{23/32})=X^{o(1)},\\
\mathcal T_0\le\exp(C\log X/\log\log X)=X^{o(1)}.
\end{gathered}
\tag{3.79}
\]

For $\alpha\in\{1-\vartheta_{\rm sp},1/2-\vartheta_{\rm sp}\}$ we have
$\sum_{p\mid a}p^{-\alpha}\le
\sum_{j\le\omega(a)}j^{-\alpha}\ll(\log X)^{1-\alpha}$.
The bound for the products then follows from
$-\log(1-v)\ll v$. The divisor bound is elementary.

To apply (3.72) to (3.68), use the cap
$H=\min(2H_0,H_*)$. Empty blocks have zero coefficient.
Increase $X_0$ in (3.6)--(3.7). Then

\[
H_2/S\ll X^{-9\sigma+\eta}L^{C_{\rm box}},\qquad H_*\le S.
\tag{3.80}
\]

The strict power margin in (3.6)--(3.7)
absorbs the fixed logarithmic factor and the factor 40.
For $S=1/2$ this gives
$H_*=40H_2\ll X^{-9\sigma+\eta}L^{C_{\rm box}}\to0$.
If the middle factor is supported at $s=1\in(1/2,1]$,
set $S=1/2$. Then $H_*<1$, every nonzero-frequency block is empty,
and $E_2^*=0$. Thus Theorem 2.8 is needed only for $S\ge1$;
in that case a nonempty block has $H\ge1$.
Each summand of (3.68) has modulus at most 1, so
$|\beta_{z,m}|\le B_{z,\varepsilon_y\varsigma_fm}$.
Since $p_0\nmid a$, $(p_0m)_a=m_a$ exactly. Thus

\[
\sum_{n''\mid a^\infty}(n'')^{\vartheta_{\rm sp}}
                 \|\beta^\flat(n'')\|_2^2
 \ll X^{o(1)}(H_*S)^2F,
\qquad\|\beta\|_2^2\ll X^{o(1)}(H_*S)^2F.
\tag{3.81}
\]

The second estimate is the same proof with $\vartheta_{\rm sp}=0$.
Here $S\ge1$, as (3.72) requires; the case $S=1/2$ was treated above.
If $\gamma_1$ is supported at $1$, take $q=1,Q=1$.
If $\gamma_3$ is supported at $1$, it is a fixed coordinate.
If the full modulus is 1, the discrepancy is identically zero.

\subsection{Application of Theorem 2.8 and the spectral budget}\label{sec:pi-spectral}

For nonzero (3.69), set

\[
\begin{aligned}
&r_T=z\asymp Z,\ s_T=p_0,\ R_T=Z,\ S_T=p_0/2,\quad\\
&d_T=b,\ c_T=c,\ n_T=p_0m,\ N_T=8p_0Y,
\quad a_T=\varepsilon_y a.
\end{aligned}
\tag{3.82}
\]

Let the Theorem 2.8 coefficient equal $\beta_{z,m}$ at these
indices and zero elsewhere. The support conditions give exactly

\[
(zb,p_0c)=1\iff
(b,c)=(z,c)=(z,p_0)=(b,p_0)=1,
\qquad(a_T,zp_0)=1.
\tag{3.83}
\]

Here $(z,c)=1$ comes from $(s,ac)=(s',ac)=1$.
The $p_0$ conditions on $(b,z)$ come from their roughness.
The smoothed variable $c$ may be divisible by $p_0$.
Reducing the inverse modulo $c$ gives the exact phase identity

\[
e\left(\frac{a_T(p_0m)\overline{zb}_{p_0c}}{p_0c}\right)
 =e\left(\frac{\varepsilon_y a m\overline{zb}_c}{c}\right).
\tag{3.84}
\]

We make the dyadic support of the smooth weight explicit.
Choose a fixed smooth step $\chi$. It is zero on $(-\infty,0]$ and one
on $[1,\infty)$. Split

\[
\begin{aligned}
\psi_{0,1}(v)&=\psi_0(v)[1-\chi((v-0.8)/0.1)],\\
\psi_{0,2}(v)&=\psi_0(v)\chi((v-0.8)/0.1)[1-\chi((v-1.4)/0.2)],\\
\psi_{0,3}(v)&=\psi_0(v)\chi((v-1.4)/0.2).
\end{aligned}
\tag{3.85}
\]

They sum to $\psi_0$.
They lie in the strict dyadic windows $(0.45,0.90]$, $(0.80,1.60]$, $(1.40,2.80]$.
Expanding the $(b,c)$ bumps gives nine weights.
Multiply by fixed smooth cutoffs equal to 1 on the supported $n_T,z,p_0$ indices.
At the lower $n_T$ endpoint use a cutoff zero below $1/2$ and one above 1.
Its derivatives are $O_\nu(n_T^{-\nu})$.
Each weight satisfies the Theorem 2.8 derivative conditions with
$\epsilon_0=\sigma/2000>0$. Take $\epsilon=\sigma/100$,
$A_T=A_0$. Thus (3.14), exactly as stated, applies to (3.69).

We bound the four terms of its bracket. Put

\[
\begin{gathered}
V=T^4R^4S^6,\quad D_*=TR,\quad C_*=TRS,\\
Z\asymp S^2,\quad Y\ll X^{o(1)+\eta}T^2R^2S^3/X.
\end{gathered}
\tag{3.86}
\]

Since $\vartheta_{\rm sp}<1/2$, all powers of $(D,C)$ in (3.15) are nonnegative.
First replace them monotonically by $D_*,C_*$.
Then $C_*\le D_*Z$, and $Y\ll ZX^{-8\sigma+\eta+o(1)}$.
Apply (3.81). Suppress only uniform $X^{o(1)}$ and $p_0$-powers. This gives

\[
\begin{array}{c|c}
\text{bracket term}/V&\text{upper bound}\\ \hline
1&K_A=A_0^{\vartheta_{\rm sp}} TQR^2S^5/X^2\\
2&(A_0/X)K_A\\
3&K_R=(TR/S)^{2\vartheta_{\rm sp}}TQR^2S^5/X^2\\
4&(TQR^2S^5/X^2)(T^2R^2/X).
\end{array}
\tag{3.87}
\]

We check the regular third term directly.
Up to $p_0$-powers, its bracket before the coefficient norm is

\[
\begin{aligned}
C(CD\sqrt Z)^{2\vartheta_{\rm sp}}Z^{1-\vartheta_{\rm sp}}
  (DZ)^{1-2\vartheta_{\rm sp}}
 &=CDZ^2(C/Z)^{2\vartheta_{\rm sp}}\\
 &\le T^2R^2S^5(TR/S)^{2\vartheta_{\rm sp}}.
\end{aligned}
\tag{3.88}
\]

The factor here is $C/Z$; the factor $C/(SZ)$ would lose $S^{2\vartheta_{\rm sp}}$.
The bound $C\le DZ$ is applied after passing to the envelopes.
The first term is at most $A_0^{\vartheta_{\rm sp}} CDZ^2$ times the weighted norm.
The second uses $A_0YZ$. The fourth uses $D^2YZ$.
Substitute $(H_*S)^2F\ll X^{o(1)+2\eta}QT^3R^4S^6/X^2$.
This proves the other lines of (3.87).

The four packing inequalities (3.7) bound the exponents in (3.87) by
$-15\sigma,-(15-1/20)\sigma,-15\sigma,-23\sigma$.
The cutoff contributes at most $3\eta$ to $J_*^2/V$.
The product in (3.14) is eventually at most $X^{10.01}$.
This follows from $A_0<X^{1.01},T<X^{0.34},R,S<X^{0.6}$ and
the sharper $Y$ bound. The bounds on $R$ and $S$ use only
$RS\le X^{0.6}$. This comes from the full-modulus support $q\le X^{0.6}$
in (H3), since the full modulus has scale $QRS$; dyadic constants are
absorbed after enlarging $X_0$.
So the smoothness prefactor costs at most $0.19019\sigma$ in the exponent.

Before the cutoff loss, the exponents relative to
$V^{1/2}$ (from the square-root row on, relative to $V^{1/4}$), in
units of $\sigma$, are

\[
\begin{array}{c|r}
J_*&-299/40=-7.475\\
\text{Theorem 2.8 prefactor}&+0.19019\\
W_{k\ne0}&-7.28481\\
W_{k=0}&-8\\
\text{square root in (3.64)}&-3.642405\\
\text{Mellin integral}&+0.075\\
\text{combined fixed subpolynomial factors}&+0.11\\
\text{cutoff }X^\eta\text{ after both roots}&+0.00075.
\end{array}
\tag{3.89}
\]

The Mellin allowance follows from (3.57).
Any fixed logarithmic power is eventually at most $X^{0.074\sigma}$.
The remaining subpolynomial factors are the finitely many $(h,f,D,C,Z)$
localizations, the divisor representation counts in (3.65), (3.71), the
prime products (3.79), and powers of $p_0\ll L$.
Their product is uniformly $X^{o(1)}$.
So it is at most $X^{0.11\sigma}$ after choosing $X_0$.
The multiplier from the random partition is a
fixed $\sigma$-dependent constant. By
(3.58), (3.64), (3.65), (3.89),

\[
|E_2^*|\ll L^C\frac{T^2}{Q}X^{-3.45\sigma},
\tag{3.90}
\]

since $-3.642405+0.075+0.11+0.00075=-3.456655<-3.45$.
For $S=1/2$, $E_2^*=0$ (see above).
Inserting (3.90) in (3.51) gives a spectral contribution
$M_ITX^{-1.725\sigma}L^C\ll X^{1-1.725\sigma}L^C$, since $M_IT\ll X$.
As $c_\sigma<\sigma/2<1.725\sigma$, it is absorbed in $X^{1-c_\sigma}L^C$.

For a target $A'$, use $A_{\rm work}=A'+4$ on every shell
triple, uniformly before summing.
The triangle inequality over $O(L^3)$ triples costs
the logarithm of the outer modulus and the
two logarithms of the other factors.
Finally restore the factor $C_{\rm slot}$ of (3.5). It is $O_\sigma(1)$.
Equations (3.51), (3.90) prove Theorem~\ref{thm:main}.

\subsection{Global summation, prime powers, and normalization}\label{sec:pi-global}

At fixed $\sigma$, assume (H7): the box family has the
source and Heath--Brown coefficient bound

\[
\sum_j|c_j|\le C_{\sigma,C_{\rm box},\mathcal C_{\rm src}}(\log N)^{K_0},\qquad
K_0\le24C_{\rm box}+26.
\tag{3.91}
\]

\textbf{Source-wide error bounds.}
Assume the source's aggregate logarithm-freezing bound and the
boundary bound (3.24).
With the exponents $c_{\log}$ and $c_R$ of (H6), a sufficient pair of
raw bounds is
\[
 E_{\rm log}\ll_\sigma N(\log N)^{c_{\log}-C_{\rm box}},\qquad
 E_{\rm bdry}\ll_\sigma N(\log N)^{c_R-C_{\rm box}/2}.
\]
They include all source labels and representation multiplicities.
These bounds are separate from the interior-box hypotheses (H1)--(H5).
The source-wide estimates enter through (B.29) for the prime source and
through the auxiliary bound (3.95). They are not hypotheses of
Theorem~\ref{thm:main}.

The source labels and prefix cells are $O_\sigma(1)$. There are
at most 24 HB variables. Hence there are $O(L^{24})$ dyadic tuples, and
each dyadic tuple contains $O(L^{24C_{\rm box}})$ short boxes.
The product shells number $O(\log\log N)\le O(L)$.
The total count is $O_\sigma(L^{24C_{\rm box}+25})$. Since the
coefficients are $O_\sigma(1)$, (3.91) holds with
$K_0=24C_{\rm box}+26$; the count exponent is smaller.
The three factor localizations are inside (3.13).
The count includes every partition.

Choose parameters in the following order.
Fix $\sigma$ and the source combinatorics.
For Theorem~\ref{thm:aa37-interior}, through (3.92) and (B.29), use
the requested exponent $A_{\rm agg}=A$ and choose
$C_{\rm box}\ge\lceil2A+32\rceil$; then set
$A'=A+24C_{\rm box}+31$ and $A_{\rm work}=A'+4$.
This is the parameter convention used in Theorem~\ref{thm:aa37-interior}
and Theorem~\ref{thm:aa37-fixed-u}.

For the stronger source-wide bookkeeping in (3.95), which is outside
Theorem~\ref{thm:main}, fix a separate final exponent $A>0$ and put
$A_{\rm agg}=A+3$.
Choose the box width $C_{\rm box}=C_{\rm box}(A_{\rm agg})$ with
\[
 C_{\rm box}(A_{\rm agg})\ge
 \max\{2(A_{\rm agg}+c_R),\ A_{\rm agg}+c_{\log}\},
\]
so that the boundary and logarithm-freezing errors are
$O(N(\log N)^{-A_{\rm agg}})$. For the Iwaniec source,
$c_R=16$ by (B.19) and the logarithm-freezing exponent is
$c_{\log}=18$ by (B.18); since
$A_{\rm agg}+18<2A_{\rm agg}+32$, the condition is the integer choice
$C_{\rm box}(A_{\rm agg})\ge\lceil2A_{\rm agg}+32\rceil
=\lceil2A+38\rceil$. Set
$A'=A_{\rm agg}+24C_{\rm box}+31$ and $A_{\rm work}=A'+4$.
The value of $A'$ is $K_0+A_{\rm agg}+5$ at the maximal count
$K_0=24C_{\rm box}+26$ in (3.91).
Let $c_*$ dominate $c_1$ and every preliminary logarithmic loss denoted by $c$.
This includes the fixed loss in (3.41) and the endpoint losses before BFI.
It is chosen before $A_E,C_\Sigma,C_M$.
Use the same $c_*$ for every source and every application below.
Choose
\[
\begin{aligned}
 A_E&\ge2A_{\rm work}+2C_{\rm box}+c_*,\\
 C_\Sigma&=A_E+2c_*+3,\\
 C_M&=5C_\Sigma+A_E+C_{\rm box}+c_*+3.
\end{aligned}
\]
Then choose $A_3$ and the required BFI/SW exponents.
Later power-saving logarithmic losses may depend on $C_\Sigma,C_M$.
The constant $c_*$ is defined before them.
Choose $N_0$ last. It is ineffective, as are the implied constants,
because the Siegel--Walfisz theorem is ineffective.

The source may choose a larger width. Its fixed width is selected
before the analytic precisions, with the same displayed formulas.

By (3.13), (H7) in the form (3.91), and $N/(\log N)^{K_{\rm scale}}\le X\le N$,

\[
\begin{aligned}
E_{\rm II}:=\sum_j|c_j||S_j^{\rm PSDT}|
&\ll N(\log N)^{K_0-A'}
 +N^{1-c_\sigma}(\log N)^{K_0+C}\\
&\ll_{\sigma,A,C_{\rm box},K_{\rm scale},\mathcal C_{\rm src}}
N(\log N)^{-A}.
\end{aligned}
\tag{3.92}
\]

Every positive fixed power saving beats every fixed logarithm once $N_0$ is chosen.
The spectral term $X^{1-1.725\sigma}L^C$ is already absorbed in (3.13), since $c_\sigma<\sigma/2$.

\textbf{Source prime powers.} For the direct source the correction in (3.22) has
$O(N^{1/2}(\log N)^C)$ integers before discrepancy weights.
For the product source assume at most four tested coordinates,
each at least $N^\alpha$, $\alpha=100/1149$, and product $e<N$.
A proper power coordinate has reciprocal mass

\[
\sum_{N^\alpha\le p^k\le N,\ k\ge2}p^{-k}
 \ll N^{-\alpha/2}+(\log N)N^{-2\alpha/3}
 \ll N^{-\alpha/2}.
\tag{3.93}
\]

For $k=2$ compare with the tail of $\sum n^{-2}$.
For $k\ge3$ sum the integer $k$-power tail.
There are $O(\log N)$ possible exponents.
Fix one coordinate at $t$.
For $1\le r\le4$, the remaining positive coordinates
have at most $(N/t)(1+\log N)^{\max(r-2,0)}$ possibilities.
For $r=1$ no coordinate remains, and the
single admissible completion is bounded by $N/t$.
For $r=2$ the remaining coordinate is counted directly.
For $r\ge3$ the reciprocal-sum argument gives the displayed logarithmic power.
Take a union bound over the at most four possible power coordinates.
For the product discrepancy write $e$ for the tested product.
Theorem~\ref{thm:aa37-interior} supplies $N\le a\le N\log N/3$.
On the product support, $e<N\le a\ll N\log N$, so
$1\le |a-e|\ll N\log N$.
Use (3.23) and $\tau(|a-e|)\ll_\epsilon N^\epsilon$ with
$\epsilon=\alpha/4$. Then the whole source-level correction obeys

\[
E_{R0}\ll N^{1-\alpha/4}(\log N)^C
       =N^{1-25/1149}(\log N)^C
       \ll_A N(\log N)^{-A}.
\tag{3.94}
\]

This bound precedes the Heath--Brown subdivision, so its box multiplicity is 1.
For a product, the bound uses the coordinate floor $N^\alpha$ assumed
immediately before (3.93).
Thus (3.94) applies to every such product source.

The bounds (H6), the logarithm-freezing estimate, (3.92) (which uses
(H7)), and (3.94), with the fixed choices above, give

\[
\begin{aligned}
&E_{\rm II}+E_{\rm bdry}+E_{\rm log}+E_{R0}
       \ll_{\sigma,A,C_{\rm box},K_{\rm scale},\mathcal C_{\rm src}}
       N(\log N)^{-A-3}=o(U_C),\\
&U_C=C(N)N/(\log N)^2,\quad C(N)\ge c>0.
\end{aligned}
\tag{3.95}
\]

The ratio is
$O_{\sigma,A,C_{\rm box},K_{\rm scale},\mathcal C_{\rm src}}
((\log N)^{-A-1})$.

For each source family, its exact source identities give the
finite sum to which these bounds apply.
The restoration, BFI, partition, Poisson and Theorem 2.8 errors
enter once, inside $E_{\rm II}$.

For any normalization $W_{\rm norm}\gg N/(\log N)^2$, (3.92)
reads explicitly

\[
\begin{aligned}
\frac{\sum_j|c_j||S_j^{\rm PSDT}|}{W_{\rm norm}}
 &\ll(\log N)^{K_0+2-A'}\\
 &\quad+N^{-c_\sigma}(\log N)^{K_0+C+2}=o(1).
\end{aligned}
\tag{3.96}
\]

Consider the original direct prime source, before HB. Its mass is not yet replaced by a logarithmic-integral (Li) mass.
Its exact prime mass $X_A=\sum_pA_D(p)$ satisfies

\[
\sum_{(p,q)=1}A_D(p)=X_A
\tag{3.97}
\]

because $p>N/(\log N)^B>q$ for $q<N^{e_Q}$ and sufficiently large $N$.
Thus the formal and natural comparators coincide on that original source.
This holds for the source itself, before the split into signed HB terms.
If a Li mass $X_{\rm Li}$ is substituted, the exact difference is
$(X_A-X_{\rm Li})/\varphi(q)$.
Its weighted sum enters once, in the main-term mass contribution,
outside the natural-comparator correction.
For the composite product source, (3.97) may fail.

\subsection{The complete non-II cover}\label{sec:pi-nonII}

These results are separate from Theorem~\ref{thm:main}.
The source-to-box identity comes from the source itself.
For a direct shell $[X,2X)$, assume the following source conditions (SH1)--(SH5).
\begin{itemize}[leftmargin=3.2em]
\item[(SH1)] The coefficient $\gamma(q)$ is independent of tested variables, $(q,2N)=1$,
$|\gamma(q)|\le C_\sigma$, $q\le N^{e_Q}$, $e_Q<153/256$,
and $(a,q)=1$ on its support.
\item[(SH2)] All tested restrictions are constant on each nonboundary box.
Smooth factor weights have $\sup|w|+V(w)\ll1$.
\item[(SH3)] All HB factors are rough at the threshold $z^*=z_0(X)$ or $z_0(N)$.
\item[(SH4)] The discrepancy is exactly (3.3).
\item[(SH5)] The assignment priority is proper prime powers,
boundary boxes, Type II, Type I-1, then Type I-2.
\end{itemize}

Put $\omega=\log(N/X)/\log N$ and
$\kappa=30\log2/\log N$. Increase $N_0$ until
$\omega,\kappa\le\sigma/202$. Let $V_i$ be the lower HB scales, and
$e_i=\log V_i/\log N$. Then (3.21) implies

\[
1-\omega-\kappa<\sum_i e_i<1+\kappa,\qquad
e_i\le1/3+\kappa\quad\text{on Möbius variables}.
\tag{3.98}
\]

There are at most six variables.
Call a variable small below $60\sigma$, and large above $1/3+14\sigma$.
Suppose no variable is in the II window.
Since $360\sigma<1-\omega-\kappa$, some variable is large.
Since $1+42\sigma>1+\kappa$, at most two are large.
Möbius variables are never large.
Thus there are exactly one or two large smooth variables.
For a window variable,

\[
\log_X(2V)\le(1/3+14\sigma+\kappa)/(1-\omega)
       \le1/3+15\sigma,
\tag{3.99}
\]

which proves the shell-local window transfer. For one large variable,

\[
e_*-e_Q-(1-\omega)/10
 \ge0.9-153/256-0.9\omega-\kappa-300\sigma>0.3023;
\tag{3.100}
\]

hence its scale $V\ge qX^{1/10}$. With two large variables,
$N_1\ge N_2\ge X^{1/3}$. The at most four other factors are small.
Their product $m$ is split dyadically, $m\asymp M$, with

\[
M\le X^{2.401\cdot10^{-8}},\qquad N_1N_2M\le2^{30}X,\qquad q<X^{0.5977}.
\tag{3.101}
\]

These follow from
$(240\sigma+\kappa)/(1-\omega)<2.401\cdot10^{-8}$,
which uses $\sigma\le10^{-10}$,
and $e_Q/(1-\omega)<0.5977$.
Consider a product source with three or four coordinates at least $N^{100/1149}$.
Its smallest coordinate has exponent between $100/1149$ and $(1+\kappa)/3$.
It has at most six HB factors.
One of them has exponent exceeding $(100/1149-\kappa)/6>60\sigma$.
Each is below $1/3+14\sigma$.
Thus, after proper prime powers and boundary boxes are removed,
all its boxes are of type II.

\textbf{Elementary sieve sandwich.} For a finite prime set $\mathcal P$
below $z$, let $P=\prod_{p\in\mathcal P}p$,
$E=\sum_{p\in\mathcal P}1/p$, and

\[
\lambda_d^{(J)}=\mu(d)1_{d\mid P,\ \omega(d)\le J}.
\tag{3.102}
\]

Let $r=\omega((n,P))$. For $r\ge1$, Pascal's identity gives
$\sum_{j\le J}(-1)^j\binom rj=(-1)^J\binom{r-1}J$. The sum is 1 for $r=0$. Hence

\[
\begin{gathered}
\Psi^-(n):=\sum_{d\mid n}\lambda_d^{(2k-1)}
 \le1_{(n,P)=1}\le \Psi^+(n):=\sum_{d\mid n}\lambda_d^{(2k)},\\
\Psi^+\ge0,\qquad \Psi^+-\Psi^-=\binom r{2k}\ge0.
\end{gathered}
\tag{3.103}
\]

The weights are 1-bounded and supported on $d\le z^{2k}$.
Also
$\sum_{d\mid P,\omega(d)=j}1/d\le E^j/j!$.
Let $s=\log D/\log z\gg(\log\log X)^3$, and choose
$2k=2\lfloor s/2\rfloor$. Mertens' bound
$E\le\log\log z+O(1)$, with $s\gg(\log\log X)^3$, gives
$2k\ge s-2\ge e^2E$ for large $X$. With Stirling,

\[
E^{2k}/(2k)!\le(eE/(2k))^{2k}\le e^{2-s}
                   \ll_A L^{-A}.
\tag{3.104}
\]

It is the sieve approximation used in I-1 and I-2.

\setcounter{lemma}{2}
\begin{lemma}[Type I, one large variable (I-1)]\label{lem:typeI1}
Assume (SH1)--(SH5).
Consider a box with one large variable.
Let the complementary coefficient satisfy $|\alpha_m|\ll\tau(m)^{B_0}$.
Uniformly for every integer residue $a$ coprime to each supported $q$,

\[
\sum_q|\gamma(q)|\left|\sum_m\alpha_m
 \sum_{v\in I,\ P^-(v)\ge z^*}w(v)\mathcal S_{mv}(q;a)\right|
 \ll_{\sigma,A,B_0}XL^{-A}.
\tag{3.105}
\]
\end{lemma}

\begin{proof}
If $(m,q)>1$, the summand vanishes. Otherwise put
$b=a\bar m(q)$.
The progression and the reduced comparator contain only integers prime to $q$.
Apply (3.102)--(3.104) to the $q$-dependent set $\mathcal P_q=\{p<z^*:p\nmid q\}$, with
$D=X^{1/20}$. For every $d\mid P_q$, $(d,q)=1$, and

\[
\#\{v\le t:v\equiv b(q),d\mid v\}=t/(dq)+O(1),
\]
\[
\#\{v\le t:(v,q)=1,d\mid v\}
 =(\varphi(q)/q)(t/d)+O(\tau(q)).
\tag{3.106}
\]

After comparator normalization, the upper and lower sandwiches have the same main term.
Their gap gives discrepancy
$\ll(t/q)e^{2-s}+2D\tau(q)$, where
$s\ge(\log\log X)^3/40$.
Endpoint subtraction and partial summation cost $O(\sup|w|+V(w))$.
By (3.100), $D\tau(q)\ll(V/q)X^{-1/20+o(1)}$.
Thus the inner discrepancy is $O_A(V/(qL^A))$.
Sum with $\sum_m|\alpha_m|\ll ML^C$, $MV\asymp X$, and
$\sum_q1/q\ll L$. Then request a larger $A$. This proves
(3.105). The sieve weights use the prime set $\mathcal P_q$.
\end{proof}

\begin{lemma}[Completed hyperbola]\label{lem:hyperbola}
Let $c\bmod q$ be a unit. For integer
intervals $J_1,J_2$, put

\[
\mathcal D=\#\{(u,v)\in J_1\times J_2:uv\equiv c(q)\}
       -\frac{\#J_1^*\#J_2^*}{\varphi(q)}.
\tag{3.107}
\]

The full-modulus Kloosterman-sum bound
\cite[Corollary~11.12]{IK04},

\[
|S(h,k;q)|\le\tau(q)q^{1/2}(h,k,q)^{1/2},
\tag{3.108}
\]
implies, uniformly in interval lengths and positions,

\[
|\mathcal D|\ll\tau(q)q^{1/2}\log^2(2q)
        +\tau(q)^2\log^2(2q)/\varphi(q).
\tag{3.109}
\]
\end{lemma}

\begin{proof}
Put
\[
F_i(h)=\sum_{u\in J_i}e(hu/q).
\]
Additive completion then gives
\[
\mathcal D=q^{-2}\sum_{h,k}F_1(-h)F_2(-k)
\left(S(h,kc;q)-\frac{c_q(h)c_q(k)}{\varphi(q)}\right).
\]

The bracket vanishes if
$h=0$ or $k=0$, so complete periods cancel. For $0<h<q$,
$|F_i(h)|\le q/(2\min(h,q-h))$, and

\[
\sum_{0<h<q,d\mid h}|F_i(h)|\ll(q/d)\log(2q).
\tag{3.110}
\]

Use $(h,k,q)^{1/2}\le\sum_{d\mid(h,k,q)}d^{1/2}$ in (3.108)
and sum $\sum d^{-3/2}$. For the principal part use
$|c_q(h)|\le(h,q)=\sum_{d\mid(h,q)}\varphi(d)$, so its
one-variable weighted sum is $O(q\tau(q)\log(2q))$. This proves
(3.109). Modulo 1 the discrepancy is zero.
\end{proof}

\begin{lemma}[Type I, two large variables (I-2)]\label{lem:typeI2}
Assume (SH1)--(SH5) and (3.101).
Consider two large smooth variables with nonnegative weights $w_1,w_2$.
Let $|\alpha_m|\ll\tau(m)^{B_0}$.
Uniformly for every integer residue $a$ coprime to each supported $q$,

\[
\begin{aligned}
&\sum_{q\asymp Q}|\gamma(q)|\Bigl|\sum_m\alpha_m
 \sum_{n_i\in I_i,\ P^-(n_i)\ge z^*}
        w_1(n_1)w_2(n_2)\mathcal S_{mn_1n_2}(q;a)\Bigr|\\
&\qquad\ll_A XL^{-A}.
\end{aligned}
\tag{3.111}
\]
\end{lemma}

The smooth factor weights in the prime-source application of
Appendix~\ref{app:proof-aa37} are nonnegative.
For a general real weight, write $w=(w+c)-c$,
with $c=\sup|w|$, and apply the lemma to the nonnegative pieces.

\begin{proof}[Proof of Lemma~\ref{lem:typeI2}]
Again restrict to $(m,q)=1$. Let $C(FG)$ and $M_q(FG)$
be the positive congruence and comparator functionals with the
positive weights $w_1,w_2$. Put $\Delta=C-M_q$.
Apply (3.103) to all primes below $z^*$, with
$D_f=X^{\epsilon_I/2}$, $\epsilon_I=\delta/4=7.5\sigma$.
Include the Brun--Bonferroni threshold for this smaller level in $N_0$.
For every fixed logarithmic exponent $A_{\rm sv}$, put
\[
s_f=\frac{\log D_f}{\log z^*},\qquad
E=\sum_{p<z^*}\frac1p,\qquad J=2\lfloor s_f/2\rfloor,
\]
and require $J\ge\max\{2,e^2E,A_{\rm sv}\log L\}$.
At fixed $\sigma>0$ this holds eventually,
since $s_f\asymp\sigma(\log\log N)^3$ and $E=O(\log\log N)$.
Hence $(z^*)^J\le D_f$, $J-1\ge1$, and
\[
\frac{E^J}{J!}\le\left(\frac{eE}{J}\right)^J\le L^{-A_{\rm sv}}.
\]
Choose $A_{\rm sv}$ after the fixed logarithmic losses in (3.115).
So the same final choice of $N_0$ covers the use of (3.103)--(3.104) in I-2.
Write $R_i=1_{P^-(n_i)\ge z^*}$, $G_i=\Psi_i^+-\Psi_i^-\ge0$.
Since $0\le R_i\le \Psi_i^+$,

\[
\Psi_1^+\Psi_2^+-G_1\Psi_2^+-\Psi_1^+G_2\le R_1R_2\le \Psi_1^+\Psi_2^+.
\tag{3.112}
\]

Since $C,M_q$ are positive,

\[
\begin{aligned}
|\Delta(R_1R_2)|\le{}&|\Delta(\Psi_1^+\Psi_2^+)|+|\Delta(G_1\Psi_2^+)|
 +|\Delta(\Psi_1^+G_2)|\\
&+M_q(G_1\Psi_2^+)+M_q(\Psi_1^+G_2).
\end{aligned}
\tag{3.113}
\]

The lower weights $\Psi_i^-$ may be negative.
The bound (3.113) contains only the nonnegative $\Psi_i^+,G_i$.
Expand the first three terms over $d_i\le D_f$.
Their coefficients are bounded by 1. Terms with $(md_1d_2,q)>1$ vanish.
Otherwise put $c\equiv a\overline{md_1d_2}\pmod q$.
This is a unit, and $u_i=n_i/d_i$ lies on the modular hyperbola
$u_1u_2\equiv c\pmod q$.
Two-variable partial summation and (3.109) give a total
$\ll_\epsilon D_f^2q^{1/2}X^\epsilon$.

For the two main gaps, positivity lets us drop the restriction
$(n_1n_2,q)=1$. Full-integer progression counts and (3.104) give

\[
\sum_{n\in I_i}G_i(n)w_i(n)\ll_{A_{\rm sv}} N_iL^{-A_{\rm sv}}+D_f,
\quad\sum_{n\in I_i}\Psi_i^+(n)w_i(n)\ll N_i+D_f.
\tag{3.114}
\]

The progression counts range over all integers.
Summing over $(m,q)$ now gives

\[
\ll M D_f^2Q^{3/2}X^\epsilon L^C+XL^{C-A_{\rm sv}}
       +M D_f(N_1+N_2)L^C+M D_f^2L^C.
\tag{3.115}
\]

By (3.101), the first exponent is less than
$2.401\cdot10^{-8}+7.5\cdot10^{-10}+1.5(0.5977)+\epsilon<0.897$
for a sufficiently small fixed $\epsilon>0$. Also
$M(N_1+N_2)\le2^{31}X/N_2\ll X^{2/3}$.
So the exponents of the last gap terms are less than 0.668.
Choose $A_{\rm sv}$ after the fixed log losses.
This proves (3.111). Squareful moduli are covered by (3.108).
\end{proof}

Finally there are $O_\sigma(L^K)$ non-II direct boxes and shell
pieces for a fixed $K$. Request (3.105), (3.111) with exponents
larger than $K+A+5$. Their sum is then $O(NL^{-A})$; in particular,
divided by $N/(\log N)^2$ it tends to zero. For the prime source these
are the exponents $A_i=K_I+A+5$ of Appendix~\ref{app:proof-aa37}.
The I-1 and I-2 boxes are disjoint parts of the non-II contribution.
Under its coordinate hypotheses, the product source has none.
Their contributions lie outside $E_{\rm II}$.

For the prime source, Section~4 supplies the packing construction and
Proposition~\ref{prop:aa37-prime-source} in Appendix~\ref{app:proof-aa37}
gives the prime-source estimate.

